% =========================================================
%  main.tex
%  BUKU AJAR PENGANTAR PEMROGRAMAN GAME (MWCA009)
%  Master document
%
%  Program Studi Teknik Informatika
%  Fakultas Teknologi Informasi
%  Universitas Bale Bandung
%
%  Kompilasi seluruh buku:
%    latexmk -pdf main.tex
%  atau:
%    pdflatex main.tex   (jalankan dua kali agar daftar isi terisi)
%
%  Kompilasi satu bab:
%    cd chapters\05
%    latexmk -pdf chapter.tex
% =========================================================

\documentclass[12pt,a4paper,oneside]{book}

% =========================================================
%  KONFIGURASI GLOBAL
% =========================================================

% =========================================================
%  preamble.tex
%  Penghubung konfigurasi global buku ajar.
%  Konten buku tidak boleh bercampur dengan konfigurasi teknis.
% =========================================================

% =========================================================
%  config/packages.tex
%  Seluruh paket LaTeX untuk Buku Ajar Pengantar Pemrograman Game
%  Aturan: paket global hanya didefinisikan di berkas ini,
%  tidak di chapter.tex atau section-*.tex (Arsitektur §36, Rule 3).
% =========================================================

% --- Pengkodean & bahasa --------------------------------
\usepackage[utf8]{inputenc}
\usepackage[T1]{fontenc}
\usepackage{lmodern}
\usepackage[english,indonesian]{babel}

% --- Tata letak halaman ---------------------------------
\usepackage[
  a4paper,
  left=2.9cm, right=2.4cm, top=2.8cm, bottom=2.8cm,
  headsep=0.5cm, footskip=1.0cm
]{geometry}
\usepackage{pdflscape}
\usepackage{fancyhdr}
\usepackage{ragged2e}
\usepackage{microtype}

% --- Tabel ----------------------------------------------
% Catatan: opsi [table] pada xcolor (colortbl) tidak dipakai karena
% colortbl pada instalasi ini menuntut array.sty yang lebih baru
% (\insert@pcolumn tidak terdefinisi). Pewarnaan baris tabel digantikan
% oleh penebalan garis header pada config/styles.tex.
\usepackage{booktabs}
\usepackage{array}
\usepackage{tabularx}
\usepackage{longtable}
\usepackage{multirow}
\usepackage{makecell}
\usepackage{xcolor}

% --- Gambar & posisi mengambang -------------------------
\usepackage{graphicx}
\usepackage{float}

% --- Tipografi & daftar ---------------------------------
\usepackage{enumitem}
\usepackage{titlesec}
\usepackage{tocloft}
\usepackage{caption}
\usepackage{needspace}
\usepackage{etoolbox}

% --- Matematika -----------------------------------------
\usepackage{amsmath}
\usepackage{amssymb}
\usepackage{amsthm}
\usepackage{mathtools}

% --- Kode program ---------------------------------------
\usepackage{listings}
\usepackage{upquote}

% --- Diagram --------------------------------------------
\usepackage{tikz}
\usetikzlibrary{
  arrows.meta,
  positioning,
  shapes.geometric,
  shapes.misc,
  calc,
  fit,
  backgrounds,
  matrix,
  chains,
  decorations.pathreplacing,
  patterns
}
\usepackage{pgfplots}
\pgfplotsset{compat=1.18}

% --- Kotak informasi ------------------------------------
\usepackage{tcolorbox}

% --- Kompilasi modular ----------------------------------
% subfiles memuat paket import: ketika \subfile{chapters/XX/chapter}
% dijalankan, \input@path dan \Ginput@path disetel ke direktori bab,
% sehingga \includegraphics{figures/...}, \input{contoh}, dan
% \subfile{section-01} tetap benar pada kompilasi bab mandiri maupun
% pada kompilasi seluruh buku.
\usepackage{subfiles}

% --- Rujukan silang antar-bab untuk kompilasi bab mandiri
% (Arsitektur §40). Label dari main.aux diimpor bila berkasnya ada.
% Dipakai xr-hyper (bukan xr) karena hyperref sudah dimuat; xr biasa
% menimbulkan galat "\XR@prefix undefined" pada \cite.
\usepackage{xr-hyper}

% --- Tautan (dimuat paling akhir) -----------------------
\usepackage{url}
\usepackage{xurl}
\usepackage[
  bookmarks=true,
  bookmarksnumbered=true,
  bookmarksopen=true,
  breaklinks=true
]{hyperref}

% =========================================================
%  config/styles.tex
%  Warna, tipe kolom, format judul, header/footer, keterangan tabel
%  Skema warna dan tipe kolom diambil dari rps/config/styles.tex
%  agar tampilan buku ajar dan RPS satu keluarga.
% =========================================================

% --- Warna institusi ------------------------------------
\definecolor{ubPrimary}{HTML}{1F3864}
\definecolor{ubAccent}{HTML}{2E75B6}
\definecolor{ubHead}{HTML}{D9E2F3}
\definecolor{ubAlt}{HTML}{F2F6FC}
\definecolor{ubLine}{HTML}{BFBFBF}

% --- Warna semantik untuk kotak & kode ------------------
\definecolor{ubGreen}{HTML}{1E7B45}
\definecolor{ubGreenBg}{HTML}{F1F8F2}
\definecolor{ubRed}{HTML}{B3261E}
\definecolor{ubRedBg}{HTML}{FDF2F1}
\definecolor{ubWarn}{HTML}{A35A05}
\definecolor{ubWarnBg}{HTML}{FFF8EC}
\definecolor{ubGray}{HTML}{F6F6F6}

% --- Tipe kolom -----------------------------------------
\newcolumntype{L}[1]{>{\RaggedRight\arraybackslash}p{#1}}
\newcolumntype{C}[1]{>{\Centering\arraybackslash}p{#1}}
\newcolumntype{J}[1]{>{\justifying\arraybackslash}p{#1}}
\newcolumntype{Y}{>{\RaggedRight\arraybackslash}X}

% --- Jarak antar paragraf -------------------------------
\setlength{\parindent}{0pt}
\setlength{\parskip}{0.5em}
\renewcommand{\arraystretch}{1.25}

% --- Toleransi pemenggalan baris ------------------------
\setlength{\emergencystretch}{2em}
\tolerance=1500
\hyphenpenalty=200
\widowpenalty=1000
\clubpenalty=1000

% --- Nama bagian buku (bahasa Indonesia) ----------------
% babel menetapkan nama bagian pada saat \begin{document}, sehingga
% \renewcommand biasa pada preamble akan ditimpa. Karena itu nama bagian
% diubah melalui mekanisme babel (\addto\captionsindonesian).
\addto\captionsindonesian{%
  \renewcommand{\bibname}{Daftar Pustaka}%
  \renewcommand{\chaptername}{Bab}%
  \renewcommand{\contentsname}{Daftar Isi}%
  \renewcommand{\listfigurename}{Daftar Gambar}%
  \renewcommand{\listtablename}{Daftar Tabel}%
  \renewcommand{\figurename}{Gambar}%
  \renewcommand{\tablename}{Tabel}%
  \renewcommand{\appendixname}{Lampiran}%
}

% --- Format judul bab -----------------------------------
\titleformat{\chapter}[display]
  {\normalfont\LARGE\bfseries\color{ubPrimary}}
  {\filright\normalsize\bfseries\color{ubAccent}\MakeUppercase{\chaptertitlename}~\thechapter}
  {0.7ex}
  {\Large}
  [\vspace{0.5ex}{\color{ubAccent}\titlerule[0.9pt]}]
\titlespacing*{\chapter}{0pt}{6pt}{26pt}

% --- Format judul section -------------------------------
\titleformat{\section}
  {\normalfont\large\bfseries\color{ubPrimary}}
  {\thesection}{0.5em}{}
\titlespacing*{\section}{0pt}{1.5em}{0.6em}

\titleformat{\subsection}
  {\normalfont\normalsize\bfseries\color{ubAccent}}
  {\thesubsection}{0.5em}{}
\titlespacing*{\subsection}{0pt}{1.1em}{0.4em}

\titleformat{\subsubsection}
  {\normalfont\normalsize\bfseries}
  {\thesubsubsection}{0.5em}{}
\titlespacing*{\subsubsection}{0pt}{0.9em}{0.3em}

% --- Format judul bab tak bernomor ----------------------
\titleformat{name=\chapter,numberless}[display]
  {\normalfont\LARGE\bfseries\color{ubPrimary}}
  {}
  {0pt}
  {\Large}
  [\vspace{0.5ex}{\color{ubAccent}\titlerule[0.9pt]}]
\titlespacing*{name=\chapter,numberless}{0pt}{6pt}{26pt}

% --- Header & footer ------------------------------------
% Front matter memakai gaya polos (nomor halaman di tengah bawah);
% gaya fancy dengan kepala halaman dipakai mulai main matter.
\pagestyle{plain}
\fancypagestyle{mainmatter}{%
  \fancyhf{}%
  \fancyhead[L]{\footnotesize\color{ubPrimary}\NamaPT}%
  \fancyhead[R]{\footnotesize\color{ubPrimary}\leftmark}%
  \fancyfoot[C]{\footnotesize\thepage}%
  \renewcommand{\headrulewidth}{0.4pt}%
  \renewcommand{\footrulewidth}{0pt}%
}
% \@chapapp bernilai "Bab" pada bagian utama dan berubah menjadi "Lampiran"
% setelah \appendix. Memakai \chaptername secara langsung akan membuat judul
% berjalan pada lampiran tertulis "Bab B", sebab \chaptername tidak berubah.
\makeatletter
\renewcommand{\chaptermark}[1]{\markboth{\@chapapp\ \thechapter.\ #1}{}}
\makeatother
\renewcommand{\sectionmark}[1]{\markright{\thesection\ #1}}

% --- Keterangan tabel & gambar --------------------------
\captionsetup{
  font=small,
  labelfont={bf,color=ubPrimary},
  skip=4pt
}

% --- Daftar (list) --------------------------------------
\setlist[itemize]{topsep=2pt,itemsep=1pt,parsep=0pt,leftmargin=1.3em}
\setlist[enumerate]{topsep=2pt,itemsep=2pt,parsep=0pt,leftmargin=1.5em}

% Daftar khusus untuk checklist kompetensi
\newlist{checklist}{itemize}{1}
\setlist[checklist]{label=$\square$,leftmargin=1.7em,itemsep=3pt,topsep=3pt}

% Daftar khusus untuk langkah praktikum
\newlist{langkah}{enumerate}{1}
\setlist[langkah]{label=\textbf{\arabic*.},leftmargin=1.9em,itemsep=3pt,topsep=3pt}

% --- Tabel ringkas (keterangan di bawah tabel) ----------
\newcommand{\tabsource}[1]{%
  \par\vspace{2pt}{\footnotesize\itshape #1\par}%
}

% --- Garis tabel (booktabs) -----------------------------
% Header tabel ditandai garis tebal di atas dan tipis di bawah,
% sebagai pengganti pewarnaan latar baris (colortbl tidak tersedia).
\renewcommand{\toprule}{\specialrule{1.3pt}{0pt}{0pt}}
\renewcommand{\midrule}{\specialrule{0.5pt}{0pt}{0pt}}

% --- Daftar isi (tocloft) -------------------------------
\renewcommand{\cftchapfont}{\bfseries\color{ubPrimary}}
\renewcommand{\cftchappagefont}{\bfseries\color{ubPrimary}}
\renewcommand{\cftsecfont}{\color{black}}
\renewcommand{\cfttoctitlefont}{\hfill\Large\bfseries\color{ubPrimary}}
\renewcommand{\cftaftertoctitle}{\hfill}
\renewcommand{\cftloftitlefont}{\hfill\Large\bfseries\color{ubPrimary}}
\renewcommand{\cftafterloftitle}{\hfill}
\renewcommand{\cftlottitlefont}{\hfill\Large\bfseries\color{ubPrimary}}
\renewcommand{\cftafterlottitle}{\hfill}

% --- Tampilan tautan ------------------------------------
\hypersetup{
  colorlinks = true,
  linkcolor  = ubPrimary,
  urlcolor   = ubAccent,
  citecolor  = ubAccent,
  filecolor  = ubAccent
}

% =========================================================
%  config/commands.tex
%  Perintah khusus buku ajar.
%  Tujuan: mencegah definisi berulang yang tersebar di chapters/.
%  Perintah \CPL, \CPMK, \SUBCPMK, \code, \istilah, \headrow,
%  \Utama, \Pendukung, \TidakAda, \tabsource diambil dari
%  rps/config/commands.tex dan rps/config/styles.tex agar
%  penulisan buku ajar dan RPS konsisten.
% =========================================================

% --- Kode capaian pembelajaran --------------------------
\newcommand{\CPL}[1]{\textbf{CPL#1}}
\newcommand{\CPMK}[1]{\textbf{CPMK#1}}
\newcommand{\SUBCPMK}[1]{\textbf{Sub-CPMK#1}}

% --- Istilah teknis -------------------------------------
% \istilah  : istilah asing (miring)
% \code     : potongan kode, nama berkas, atau nama fungsi di dalam teks
%
% \code memakai \detokenize sehingga garis bawah dan karakter khusus lain
% boleh ditulis APA ADANYA tanpa perlu di-escape:
%     benar : \code{pygame.key.get_pressed}   \code{code/python/uji_gerak.py}
%     salah : \code{pygame.key.get\_pressed}  (akan tercetak "\_")
\newcommand{\code}[1]{\texttt{\detokenize{#1}}}
\newcommand{\istilah}[1]{\textit{#1}}

% --- Penanda baris header tabel -------------------------
\newcommand{\headrow}{}

% --- Penanda status hubungan (matriks CPL--CPMK) --------
\newcommand{\Utama}{\textbf{Utama}}
\newcommand{\Pendukung}{Pendukung}
\newcommand{\TidakAda}{---}

% --- Penanda indikator pencapaian -----------------------
% Dipakai pada tabel Sub-CPMK: \ind{a}, \ind{b}, \ind{c}
\newcommand{\ind}[1]{\textbf{(#1)}}

% --- Keterangan sumber gambar ---------------------------
% Contoh: \sumber{[R12] hal. 24}
\newcommand{\sumber}[1]{%
  \par\vspace{2pt}{\footnotesize\itshape Sumber: #1\par}%
}

% --- Akar direktori yang bergantung pada cara kompilasi ----
% Ketika sebuah bab dikompilasi mandiri (misalnya latexmk -pdf chapter.tex
% di dalam chapters/05), direktori kerjanya adalah direktori bab sehingga
% akar buku berada dua tingkat di atas. Ketika seluruh buku dikompilasi,
% direktori kerjanya adalah akar buku.
%
% Kondisi \ifSubfilesClassLoaded dievaluasi SEKALI di sini, yaitu ketika
% berkas ini dibaca sebagai bagian dari preamble.
%
% PENTING: \ifSubfilesClassLoaded BUKAN kondisional gaya \if...\else...\fi.
% Perintah ini menerima DUA argumen dalam kurung kurawal, seperti
% \@ifundefined: \ifSubfilesClassLoaded{<bila mandiri>}{<bila seluruh buku>}.
% Menulisnya sebagai \if...\else...\fi menyebabkan galat "Extra \else" dan
% menghasilkan nama berkas yang kacau.
%
% Kondisi juga tidak boleh diletakkan di dalam argumen nama berkas
% (\includegraphics, \lstinputlisting) karena argumen tersebut diproses
% ulang oleh \edef.
\ifSubfilesClassLoaded{%
  \newcommand{\coderoot}{../../code/}%
  \newcommand{\assetroot}{../../assets/}%
}{%
  \newcommand{\coderoot}{code/}%
  \newcommand{\assetroot}{assets/}%
}

% Dipakai sebagai: \lstinputlisting{\codefile{python/ch04_game_loop.py}}
\newcommand{\codefile}[1]{\coderoot#1}

% Dipakai sebagai: \includegraphics{\assetfile{images/tangkapan.png}}
\newcommand{\assetfile}[1]{\assetroot#1}

% --- Rujukan ke bab lain ---------------------------------
% \refbab{chap:game-loop}{Bab tentang Game Loop}
% Mencetak nomor bab bila label tersedia (kompilasi seluruh buku atau
% label diimpor melalui paket xr-hyper); bila tidak, mencetak teks
% cadangan sehingga bab tetap dapat dikompilasi mandiri (Arsitektur §40).
\makeatletter
\newcommand{\refbab}[2]{%
  \@ifundefined{r@#1}{#2}{Bab~\ref{#1}}%
}
\makeatother

% --- Penanda bagian yang belum selesai (hanya untuk penyusunan) --
% Tidak boleh dipakai pada versi terbit.
\newcommand{\todo}[1]{%
  \begin{tcolorbox}[
    breakable, sharp corners, boxrule=0.5pt,
    colframe=ubRed, colback=ubRedBg, coltitle=white,
    fonttitle=\bfseries\small, title={Belum selesai}]
  \footnotesize #1
  \end{tcolorbox}%
}

% =========================================================
%  config/environments.tex
%  Lingkungan (environment) khusus buku ajar berbasis OBE.
%  Setiap kotak dipakai konsisten pada seluruh bab sehingga
%  pembaca mengenali fungsi tiap bagian.
% =========================================================

\tcbuselibrary{breakable,skins,theorems}

% --- Gaya dasar kotak buku ------------------------------
\tcbset{
  bookbox/.style={
    breakable,
    sharp corners,
    boxrule=0.7pt,
    left=7pt, right=7pt, top=5pt, bottom=5pt,
    fonttitle=\bfseries\small,
    coltitle=white,
    titlerule=0.4pt,
    toptitle=2pt, bottomtitle=2pt,
    before skip=9pt, after skip=9pt,
    before upper={\setlength{\parskip}{0.35em}}
  }
}

% --- Tujuan pembelajaran bab ----------------------------
\newtcolorbox{tujuanbox}[1][]{%
  bookbox, colframe=ubPrimary, colback=ubAlt,
  title={Tujuan Pembelajaran}, #1}

% --- Sub-CPMK yang dicapai bab --------------------------
\newtcolorbox{subcpmkbox}[1][]{%
  bookbox, colframe=ubAccent, colback=ubAlt,
  title={Sub-CPMK Bab Ini}, #1}

% --- Contoh dan ilustrasi -------------------------------
\newtcolorbox{contohbox}[1][]{%
  bookbox, colframe=ubGreen, colback=ubGreenBg,
  title={Contoh}, #1}

% --- Catatan penjelas -----------------------------------
\newtcolorbox{catatabox}[1][]{%
  bookbox, colframe=ubLine, colback=ubGray, coltitle=ubPrimary,
  title={Catatan}, #1}

% --- Kesalahan yang sering terjadi ----------------------
\newtcolorbox{peringatanbox}[1][]{%
  bookbox, colframe=ubRed, colback=ubRedBg,
  title={Perhatian}, #1}

% --- Tip praktik ----------------------------------------
\newtcolorbox{tipsbox}[1][]{%
  bookbox, colframe=ubWarn, colback=ubWarnBg,
  title={Tip Praktik}, #1}

% --- Lembar praktikum -----------------------------------
\newtcolorbox{praktikumbox}[1][]{%
  bookbox, colframe=ubPrimary, colback=white,
  title={Praktikum}, #1}

% --- Aktivitas belajar yang direkomendasikan ------------
\newtcolorbox{aktivitasbox}[1][]{%
  bookbox, colframe=ubAccent, colback=white,
  title={Aktivitas Pembelajaran yang Direkomendasikan}, #1}

% --- Asesmen dan instrumen penilaian --------------------
\newtcolorbox{asesmenbox}[1][]{%
  bookbox, colframe=ubAccent, colback=white,
  title={Asesmen dan Instrumen Penilaian}, #1}

% --- Checklist pencapaian kompetensi --------------------
\newtcolorbox{checklistbox}[1][]{%
  bookbox, colframe=ubGreen, colback=white,
  title={Checklist Pencapaian Kompetensi}, #1}

% --- Refleksi mahasiswa ---------------------------------
\newtcolorbox{refleksibox}[1][]{%
  bookbox, colframe=ubWarn, colback=white,
  title={Refleksi}, #1}

% --- Rangkuman bab --------------------------------------
\newtcolorbox{rangkumanbox}[1][]{%
  bookbox, colframe=ubPrimary, colback=ubAlt,
  title={Rangkuman}, #1}

% --- Teorema matematis (amsthm) -------------------------
\theoremstyle{definition}
\newtheorem{definisi}{Definisi}[chapter]
\newtheorem{contoh}{Contoh}[chapter]
\newtheorem{prosedur}{Prosedur}[chapter]

\theoremstyle{plain}
\newtheorem{teorema}{Teorema}[chapter]
\newtheorem{lemma}{Lemma}[chapter]

% --- Pemanggil judul bagian tanpa nomor ----------------
% \bagian{Tujuan Pembelajaran}
\newcommand{\bagian}[1]{%
  \needspace{3\baselineskip}%
  \section*{#1}%
  \addcontentsline{toc}{section}{#1}%
}

% =========================================================
%  config/listings.tex
%  Konfigurasi penampilan kode program: Python dan C# (Unity).
%  Seluruh buku memakai konfigurasi yang sama.
% =========================================================

% --- Gaya dasar kode program ----------------------------
\lstdefinestyle{bookstyle}{
  basicstyle=\ttfamily\footnotesize,
  breaklines=true,
  breakatwhitespace=false,
  postbreak=\mbox{\textcolor{ubAccent}{$\hookrightarrow$}\space},
  frame=single,
  framerule=0.4pt,
  rulecolor=\color{ubLine},
  numbers=left,
  numberstyle=\tiny\color{ubAccent},
  numbersep=9pt,
  showstringspaces=false,
  tabsize=4,
  keepspaces=true,
  columns=fullflexible,
  captionpos=b,
  xleftmargin=1.9em,
  framexleftmargin=1.5em,
  keywordstyle=\color{ubPrimary}\bfseries,
  commentstyle=\color{ubGreen}\itshape,
  stringstyle=\color{ubRed},
  emphstyle=\color{ubWarn}\bfseries,
  extendedchars=true,
  inputencoding=utf8
}

% --- Gaya keluaran program (tanpa nomor baris) ----------
\lstdefinestyle{outputstyle}{
  style=bookstyle,
  numbers=none,
  xleftmargin=1.1em,
  framexleftmargin=0.8em,
  basicstyle=\ttfamily\footnotesize\color{black},
  backgroundcolor=\color{ubGray}
}

% --- Bahasa Python --------------------------------------
\lstdefinelanguage{PythonBook}{
  language=Python,
  morekeywords={self,True,False,None,as,with,assert,lambda,yield,global,nonlocal},
  morecomment=[l]{\#},
  morestring=[b]"
}

% --- Bahasa C# (skrip Unity) ----------------------------
\lstdefinelanguage{CSharp}{
  morekeywords={
    abstract,as,base,bool,break,byte,case,catch,char,checked,class,const,
    continue,decimal,default,delegate,do,double,else,enum,event,explicit,
    extern,false,finally,fixed,float,for,foreach,goto,if,implicit,in,int,
    interface,internal,is,lock,long,namespace,new,null,object,operator,out,
    override,params,private,protected,public,readonly,ref,return,sbyte,
    sealed,short,sizeof,stackalloc,static,string,struct,switch,this,throw,
    true,try,typeof,uint,ulong,unchecked,unsafe,ushort,using,virtual,void,
    volatile,while,var,get,set,value,async,await
  },
  sensitive=true,
  morecomment=[l]{//},
  morecomment=[s]{/*}{*/},
  morestring=[b]",
  morestring=[b]'
}

% --- Istilah khusus Unity ditandai tersendiri -----------
\newcommand{\unityapi}[1]{\textcolor{ubWarn}{\textbf{#1}}}

\lstset{style=bookstyle}

% --- Perintah ringkas pemanggil kode --------------------
% \kodepython{python/ch04_game_loop.py}{Judul}{lst:label}
\newcommand{\kodepython}[3]{%
  \lstinputlisting[style=bookstyle, language=PythonBook,
    caption={#2}, label={#3}]{\codefile{#1}}%
}
% \kodecsharp{unity/ch13_player_controller.cs}{Judul}{lst:label}
\newcommand{\kodecsharp}[3]{%
  \lstinputlisting[style=bookstyle, language=CSharp,
    caption={#2}, label={#3}]{\codefile{#1}}%
}
% \kodekeluaran{...} untuk menampilkan potongan keluaran program
\lstnewenvironment{kodekeluaran}[1][]{%
  \lstset{style=outputstyle, #1}}{}


% --- Rujukan silang antar-bab pada kompilasi bab mandiri --
% Ketika sebuah bab dikompilasi sendiri (latexmk -pdf chapter.tex di
% dalam chapters/XX), label bab lain tidak tersedia. Label tersebut
% diimpor dari main.aux melalui paket xr sehingga \refbab tetap
% menghasilkan nomor yang benar (Arsitektur §40).
\ifSubfilesClassLoaded{%
  \IfFileExists{../../main.aux}{\externaldocument{../../main}}{}%
}{}

% =========================================================
%  metadata.tex
%  Metadata buku dan identitas mata kuliah.
%  Diubah di sini akan berlaku pada seluruh dokumen.
% =========================================================

% --- Buku ------------------------------------------------
\newcommand{\BookTitle}{Pengantar Pemrograman Game}
\newcommand{\BookSubtitle}{Buku Ajar Berbasis \istilah{Outcome-Based Education} (OBE)}
\newcommand{\BookEdition}{Edisi Pertama}
\newcommand{\BookYear}{2026}
\newcommand{\BookAuthor}{Cecep Suwanda, S.Si., M.Kom.}

% --- Identitas mata kuliah -------------------------------
\newcommand{\NamaMK}{Pengantar Pemrograman Game}
\newcommand{\KodeMK}{MWCA009}
\newcommand{\BobotSKS}{3 SKS (2 SKS teori, 1 SKS praktikum)}
\newcommand{\SemesterMK}{Ganjil}
\newcommand{\TahunAkademik}{2026/2027}
\newcommand{\MataKuliahPrasyarat}{Algoritma dan Pemrograman}

% --- Identitas institusi ---------------------------------
\newcommand{\NamaProdi}{Teknik Informatika}
\newcommand{\NamaFakultas}{Teknologi Informasi}
\newcommand{\NamaPT}{Universitas Bale Bandung}
\newcommand{\NamaDosen}{Cecep Suwanda, S.Si., M.Kom.}

% --- Metadata berkas PDF ---------------------------------
\hypersetup{
  pdftitle  = {\BookTitle\ --- Buku Ajar \NamaMK},
  pdfauthor = {\BookAuthor},
  pdfsubject = {Buku ajar berbasis OBE untuk mata kuliah \NamaMK\ (\KodeMK)},
  pdfkeywords = {pemrograman game, Pygame, Unity, OBE, \NamaProdi}
}


% =========================================================
%  DOKUMEN
% =========================================================

\begin{document}

% =========================================================
%  FRONT MATTER
% =========================================================

\frontmatter

% =========================================================
%  frontmatter/cover.tex
%  Halaman judul
% =========================================================

\begin{titlepage}
\centering
\thispagestyle{empty}

{\large\bfseries\color{ubPrimary}\NamaPT\par}
\vspace{0.15cm}
{\large\color{ubPrimary}\NamaFakultas\par}
{\large\color{ubPrimary}Program Studi \NamaProdi\par}

\vspace{0.7cm}
{\color{ubAccent}\rule{\textwidth}{1.2pt}}

\vspace{1.5cm}

{\Huge\bfseries\color{ubPrimary} \BookTitle\par}

\vspace{0.5cm}
{\large\color{ubPrimary}\BookSubtitle\par}

\vspace{1.8cm}

{\large Buku Ajar Mata Kuliah\par}
\vspace{0.35cm}
{\Large\bfseries \NamaMK\par}
\vspace{0.35cm}
{\large Kode Mata Kuliah: \KodeMK\par}
\vspace{0.15cm}
{\large Bobot: \BobotSKS\par}
\vspace{0.15cm}
{\large Semester \SemesterMK\ --- Tahun Akademik \TahunAkademik\par}

\vfill

{\color{ubAccent}\rule{\textwidth}{1.2pt}}
\vspace{0.7cm}

{\large\bfseries \BookAuthor\par}
\vspace{0.5cm}

{\color{ubPrimary}\bfseries \BookEdition\ --- \BookYear\par}

\vspace{0.3cm}
{\color{ubPrimary}\NamaPT\par}

\end{titlepage}

\clearpage

% =========================================================
%  frontmatter/copyright.tex
%  Halaman hak cipta dan keterangan penggunaan
% =========================================================

\chapter*{Hak Cipta dan Keterangan Penggunaan}
\addcontentsline{toc}{chapter}{Hak Cipta dan Keterangan Penggunaan}
\thispagestyle{plain}

\vspace{0.5cm}

{\bfseries \BookTitle\par}
{\itshape \BookSubtitle\par}
\vspace{0.4cm}

Hak cipta \copyright\ \BookYear\ pada \NamaDosen, Program Studi \NamaProdi,
\NamaFakultas, \NamaPT.

\vspace{0.6cm}

{\bfseries Penerbit\par}
\NamaPT\par
\NamaFakultas\par
Program Studi \NamaProdi

\vspace{0.9cm}

{\bfseries Keterangan Penggunaan\par}

Buku ajar ini disusun untuk keperluan pembelajaran mata kuliah \NamaMK\
(\KodeMK) di lingkungan \NamaPT. Penggandaan, penyalinan, atau penyebaran
sebagian maupun seluruh isi buku untuk keperluan pendidikan nonkomersial
diperbolehkan dengan menyebutkan sumbernya. Penggunaan untuk keperluan
komersial memerlukan izin tertulis dari penulis.

\vspace{0.6cm}

{\bfseries Keterangan atas Materi Pihak Ketiga\par}

Buku ini memuat rujukan dan, pada beberapa bagian, gambar pendukung yang
berasal dari sumber pihak ketiga. Setiap gambar yang bukan karya penulis
diberi keterangan sumber di bawah gambar yang bersangkutan dengan nomor
rujukan \mbox{[R1]}--\mbox{[R35]} sesuai daftar pustaka. Seluruh merek
dagang, nama produk, dan logo yang disebutkan dalam buku ini merupakan
milik pemegang haknya masing-masing, termasuk Python, Pygame, dan Unity
yang digunakan semata-mata sebagai perkakas pembelajaran.

\vspace{0.6cm}

{\bfseries Penafian\par}

Buku ajar ini mengikuti Rencana Pembelajaran Semester (RPS) mata kuliah
\NamaMK\ yang berlaku pada Program Studi \NamaProdi\ \NamaPT. Perubahan
RPS, versi perkakas, atau tautan sumber belajar di luar kendali penulis
dapat menyebabkan sebagian isi buku perlu disesuaikan pada edisi
berikutnya. Tautan daring pada daftar pustaka ditelusuri pada September
2026 dan perlu diverifikasi ulang ketika buku digunakan.

\vspace{1.0cm}

\begin{tabular}{@{}L{2.6cm}L{9.6cm}@{}}
\textbf{Edisi} & \BookEdition \\
\textbf{Tahun} & \BookYear \\
\textbf{Penulis} & \BookAuthor \\
\textbf{Mata kuliah} & \NamaMK\ (\KodeMK), \BobotSKS \\
\end{tabular}

\clearpage

% =========================================================
%  frontmatter/preface.tex
%  Kata pengantar
% =========================================================

\chapter*{Kata Pengantar}
\addcontentsline{toc}{chapter}{Kata Pengantar}
\thispagestyle{plain}

Buku ajar ini disusun untuk mendukung perkuliahan \NamaMK\ (\KodeMK) yang
berbobot \BobotSKS\ pada Program Studi \NamaProdi, \NamaFakultas,
\NamaPT. Penyusunannya berangkat dari satu pertanyaan sederhana: bagaimana
menyajikan materi pemrograman permainan agar mahasiswa benar-benar mampu
membangun sesuatu, bukan sekadar mengenal istilah.

Mata kuliah ini dirancang dengan pendekatan \istilah{Outcome-Based
Education} (OBE). Konsekuensinya, seluruh materi, aktivitas, dan asesmen
dalam buku ini diarahkan pada capaian pembelajaran yang telah ditetapkan
lebih dahulu. Setiap bab menyebut Sub-CPMK yang dicapainya, menyediakan
praktikum yang menghasilkan artefak nyata, dan menutup dengan checklist
kompetensi serta instrumen asesmen. Dengan demikian pembaca dapat menelusuri
secara eksplisit hubungan antara apa yang dipelajari, apa yang dikerjakan,
dan apa yang dinilai.

Pilihan teknologi pada buku ini mengikuti urutan yang disengaja. Materi
pertama sampai kesembilan dikerjakan dengan Python dan Pygame tanpa bantuan
\istilah{game engine}. Mahasiswa menulis sendiri \istilah{game loop},
menghitung sendiri gerak berbasis \istilah{delta time}, menerapkan sendiri
hukum fisika sederhana, dan mengimplementasikan sendiri algoritma
kecerdasan buatan untuk NPC, termasuk \istilah{Finite State Machine} dan
pencarian jalur A*. Cara ini membuat mekanisme internal permainan terlihat
jelas dan dapat diuji. Setelah fondasi tersebut terbentuk, materi kesepuluh
sampai keempat belas memindahkan pekerjaan ke Unity, tempat mahasiswa
belajar memanfaatkan \istilah{game engine} untuk membangun level, mengimpor
aset, menata antarmuka, dan mengintegrasikan audio. Perbandingan antara
kedua pendekatan ini menjadi salah satu bahan refleksi penting pada akhir
semester.

Buku ini tidak dimaksudkan sebagai pengganti dokumentasi resmi. Python,
Pygame, dan Unity memiliki dokumentasi yang lengkap dan terus diperbarui;
buku ini berperan sebagai pemandu yang mengurutkan, menjelaskan, dan
mengaitkan materi tersebut dengan kebutuhan mata kuliah. Rujukan
\mbox{[R1]}--\mbox{[R14]} pada daftar pustaka adalah sumber utama yang
seyogianya dibaca mahasiswa, sedangkan \mbox{[R15]}--\mbox{[R35]} memuat
dokumen kurikulum, bahan ajar, dan studi kasus dari berbagai perguruan
tinggi Indonesia yang menjadi bahan pembanding penyusunan buku ini.

Umpan balik dari dosen pengampu dan mahasiswa sangat berharga bagi
perbaikan edisi berikutnya. Saran dapat disampaikan langsung kepada penulis
melalui Program Studi \NamaProdi, \NamaPT.

\vspace{0.8cm}

\begin{flushright}
Bandung, \BookYear\\
\BookAuthor
\end{flushright}

\clearpage

% =========================================================
%  frontmatter/acknowledgements.tex
%  Ucapan terima kasih
% =========================================================

\chapter*{Ucapan Terima Kasih}
\addcontentsline{toc}{chapter}{Ucapan Terima Kasih}
\thispagestyle{plain}

Buku ajar ini tidak akan selesai tanpa dukungan berbagai pihak. Penulis
menyampaikan terima kasih kepada pimpinan \NamaPT, Dekan \NamaFakultas,
dan Ketua Program Studi \NamaProdi\ atas kesempatan dan dukungan yang
diberikan dalam penyusunan bahan ajar ini.

Terima kasih juga disampaikan kepada rekan dosen di lingkungan Program
Studi \NamaProdi\ yang telah memberikan masukan atas struktur materi,
khususnya pada bagian asesmen berbasis capaian pembelajaran. Masukan
tersebut sangat membantu memastikan bahwa setiap bab benar-benar terhubung
dengan Sub-CPMK, bukan hanya memuat materi yang menarik.

Penulis juga menghargai para penulis dan pengembang yang karyanya menjadi
sumber utama buku ini, yaitu dokumentasi resmi Python, Pygame, dan Unity,
buku \istilah{Making Games with Python \& Pygame} serta \istilah{Invent Your
Own Computer Games with Python} karya Al Sweigart, tulisan Robert Nystrom
tentang \istilah{Game Loop}, penjelasan Glenn Fiedler mengenai
\istilah{timestep}, bahan kuliah Edirlei dos Santos dan Rinaldi Munir
mengenai kecerdasan buatan permainan serta algoritma A*. Karya-karya itu
membuat penulisan buku ajar ini menjadi mungkin.

Akhirnya, terima kasih kepada mahasiswa yang telah dan akan mengikuti mata
kuliah ini. Pertanyaan, kesalahan yang ditemukan, dan kesulitan yang
dihadapi ketika mengerjakan praktikum adalah bahan perbaikan yang paling
jujur bagi buku ini.

\vspace{0.8cm}

\begin{flushright}
\BookAuthor
\end{flushright}

\clearpage

% =========================================================
%  frontmatter/petunjuk-penggunaan.tex
%  Cara menggunakan buku ajar
% =========================================================

\chapter*{Petunjuk Penggunaan Buku}
\addcontentsline{toc}{chapter}{Petunjuk Penggunaan Buku}
\thispagestyle{plain}

Bagian ini menjelaskan bagaimana buku ini disusun dan bagaimana cara
menggunakannya agar kegiatan belajar berjalan efektif.

\section*{Hubungan Buku dengan Rencana Pembelajaran Semester}

Buku ini bukan dokumen yang berdiri sendiri. Isinya mengikuti Rencana
Pembelajaran Semester (RPS) mata kuliah \NamaMK\ (\KodeMK) yang berlaku di
Program Studi \NamaProdi. Setiap bab inti pada buku ini setara dengan satu
materi pembelajaran pada RPS, dan satu materi tersebut menopang satu atau
beberapa Sub-CPMK. Keterkaitan lengkap antara CPL, CPMK, Sub-CPMK, bab, dan
minggu perkuliahan disajikan pada bagian \istilah{Capaian Pembelajaran dan
Peta Kompetensi} yang menyusul sebelum Bab I.

Karena itu, urutan bab sebaiknya diikuti sebagaimana adanya. Materi
Python/Pygame pada Bab III sampai Bab XI menjadi prasyarat bagi materi
Unity pada Bab XII sampai Bab XVI, dan Bab XVII baru bermakna setelah
seluruh bab sebelumnya dikerjakan.

\section*{Struktur Setiap Bab}

Bab inti pada buku ini selalu memuat bagian berikut dengan urutan yang
tetap, sehingga pembaca dapat menemukan informasi tertentu dengan cepat.

\begin{center}
\begin{tabularx}{\textwidth}{@{}L{4.3cm} Y@{}}
\toprule
\headrow \textbf{Bagian} & \textbf{Fungsi} \\
\midrule
Sub-CPMK Bab Ini &
Rumusan kompetensi yang harus dicapai beserta indikator dan tingkat
\istilah{Taksonomi Bloom}-nya. Bagian ini menjadi dasar penilaian bab. \\
\addlinespace
Tujuan Pembelajaran &
Uraian singkat apa yang mampu dilakukan pembaca setelah menuntaskan bab. \\
\addlinespace
Pendahuluan &
Pengait dengan bab sebelumnya dan peta isi bab. \\
\addlinespace
Materi Pokok &
Pembahasan utama, terbagi menjadi beberapa subbab bernomor. \\
\addlinespace
Contoh &
Contoh kerja yang dapat diikuti langkah demi langkah. \\
\addlinespace
Praktikum &
Lembar kerja yang harus dikerjakan pada komputer, memakai kode program
yang tersedia pada lampiran daring dan repositori mata kuliah. \\
\addlinespace
Latihan &
Soal latihan dan butir refleksi untuk menguji pemahaman sendiri. \\
\addlinespace
Rangkuman &
Ringkasan butir-butir penting bab. \\
\addlinespace
Checklist Pencapaian Kompetensi &
Daftar indikator yang dapat dicentang pembaca untuk menilai kesiapan
sebelum asesmen. \\
\addlinespace
Asesmen dan Instrumen Penilaian &
Bentuk asesmen bab beserta rubrik atau kunci penilaiannya. \\
\bottomrule
\end{tabularx}
\end{center}

\section*{Kotak Informasi}

Sepanjang buku ini dipakai beberapa kotak dengan makna yang tetap.

\begin{itemize}
  \item \textbf{Tujuan Pembelajaran} dan \textbf{Sub-CPMK Bab Ini}:
        kompetensi dan indikator yang menjadi sasaran bab.
  \item \textbf{Contoh}: contoh kerja yang sebaiknya diketik ulang, bukan
        hanya dibaca.
  \item \textbf{Catatan}: penjelasan tambahan atau pengait dengan konsep
        lain.
  \item \textbf{Perhatian}: kesalahan yang sering terjadi dan perlu
        dihindari.
  \item \textbf{Tip Praktik}: saran praktis yang mempercepat pekerjaan.
  \item \textbf{Praktikum}: lembar kerja yang harus diselesaikan.
  \item \textbf{Checklist Pencapaian Kompetensi}: daftar periksa mandiri.
  \item \textbf{Refleksi}: pertanyaan untuk meninjau kembali cara belajar.
\end{itemize}

\section*{Konvensi Penulisan}

\begin{itemize}
  \item Istilah asing dicetak miring, misalnya \istilah{game loop} dan
        \istilah{collision}.
  \item Nama berkas, perintah, nama fungsi, dan potongan kode di dalam teks
        dicetak dengan huruf mesin tik, misalnya \code{pygame.time.Clock}.
  \item Kode program yang panjang disajikan sebagai daftar kode bernomor
        dan disimpan sebagai berkas terpisah pada folder \code{code/} proyek
        buku, sehingga dapat dijalankan langsung di komputer.
  \item Rujukan pustaka ditandai nomor dalam kurung siku, misalnya
        \mbox{[R3]}, dan dapat dicari pada Daftar Pustaka di bagian akhir
        buku. Nomor rujukan ini sama dengan nomor rujukan pada RPS mata
        kuliah, sehingga penelusuran antara kedua dokumen menjadi mudah.
\end{itemize}

\section*{Cara Belajar yang Disarankan}

\begin{langkah}
  \item Bacalah bagian \istilah{Sub-CPMK Bab Ini} dan \istilah{Tujuan
        Pembelajaran} lebih dahulu untuk mengetahui sasaran bab.
  \item Bacalah materi pokok secara aktif: ketik ulang setiap potongan kode
        yang penting, lalu ubah nilainya untuk melihat pengaruhnya.
  \item Kerjakan praktikum sampai program berjalan. Jika program gagal,
        bandingkan kembali dengan daftar kode pada bab dan periksa pesan
        kesalahan yang muncul.
  \item Kerjakan latihan tanpa melihat jawaban terlebih dahulu, lalu
        gunakan checklist kompetensi untuk menilai diri sendiri.
  \item Simpan setiap hasil praktikum pada repositori pribadi. Artefak
        inilah yang menjadi bukti capaian pembelajaran pada penilaian
        portofolio.
  \item Jawab butir refleksi pada akhir bab secara jujur. Bagian yang terasa
        sulit menunjukkan apa yang perlu diulang.
\end{langkah}

\begin{catatabox}
Buku ini menuntut pekerjaan pada komputer. Pastikan lingkungan
pengembangan sudah siap sebelum memulai Bab III. Langkah pemasangan
Python, Pygame, dan editor diuraikan pada bab tersebut, dan pemasangan
Unity diuraikan pada Bab XII.
\end{catatabox}

\clearpage

% =========================================================
%  frontmatter/learning-outcomes.tex
%  Capaian pembelajaran dan peta kompetensi
% =========================================================

\chapter*{Capaian Pembelajaran dan Peta Kompetensi}
\addcontentsline{toc}{chapter}{Capaian Pembelajaran dan Peta Kompetensi}
\thispagestyle{plain}

Bagian ini memuat capaian pembelajaran yang dibebankan pada mata kuliah
\NamaMK\ beserta peta keterkaitannya dengan bab-bab pada buku ini. Pembaca
yang ingin mengetahui alasan sebuah materi diajarkan, atau ingin memastikan
materi mana yang menopang suatu asesmen, dapat berpangkal dari tabel-tabel
berikut.

\section*{Capaian Pembelajaran Lulusan (CPL)}

\begin{table}[htbp]
\centering
\caption{Rumusan capaian pembelajaran lulusan yang dibebankan pada mata kuliah}
\label{tab:cpl-buku}
\begin{tabularx}{\textwidth}{@{}C{1.8cm} Y@{}}
\toprule
\headrow \textbf{Kode} & \textbf{Rumusan} \\
\midrule
\CPL{03} &
Kemampuan merancang dan menganalisis algoritma untuk menyelesaikan
permasalahan secara optimal, serta memilih dan menerapkannya pada bahasa
pemrograman tertentu. \\
\addlinespace
\CPL{04} &
Kemampuan mendesain, mengimplementasi, dan mengevaluasi solusi berbasis
\istilah{computing} yang memenuhi kebutuhan \istilah{computing} pada sebuah
disiplin program. \\
\bottomrule
\end{tabularx}
\end{table}

\section*{Capaian Pembelajaran Mata Kuliah (CPMK)}

\begin{table}[htbp]
\centering
\caption{Rumusan CPMK, CPL yang didukung, dan ranah taksonomi Bloom}
\label{tab:cpmk-buku}
\begin{tabularx}{\textwidth}{@{}C{1.7cm}Y C{3.4cm} C{1.7cm}@{}}
\toprule
\headrow \textbf{Kode} & \textbf{Rumusan CPMK} & \textbf{CPL Terkait} & \textbf{Bloom} \\
\midrule
\CPMK{1} & Merancang \istilah{game loop} dan logika fisika sederhana. &
\CPL{03} (utama), \CPL{04} (pendukung) & C2--C4 \\
\addlinespace
\CPMK{2} & Menerapkan algoritma kecerdasan buatan (AI) dasar untuk NPC. &
\CPL{03} (utama), \CPL{04} (pendukung) & C3--C5 \\
\addlinespace
\CPMK{3} & Menggunakan \istilah{game engine} (Unity) untuk membangun level. &
\CPL{04} (utama), \CPL{03} (pendukung) & C3--C4 \\
\addlinespace
\CPMK{4} & Mengintegrasikan aset audio dan visual ke dalam interaksi
\istilah{game}. &
\CPL{04} (utama), \CPL{03} (pendukung) & C3--C5 \\
\bottomrule
\end{tabularx}
\end{table}

\begin{table}[htbp]
\centering
\caption{Matriks hubungan CPL dan CPMK}
\label{tab:matriks-cpl-cpmk-buku}
\begin{tabularx}{\textwidth}{@{}L{4.4cm}@{\extracolsep{\fill}}C{2.3cm} C{2.3cm} C{2.3cm} C{2.3cm}@{}}
\toprule
\headrow & \textbf{\CPMK{1}} & \textbf{\CPMK{2}} & \textbf{\CPMK{3}} & \textbf{\CPMK{4}} \\
\headrow & \footnotesize \istilah{Game loop} \& fisika & \footnotesize AI NPC & \footnotesize \istilah{Engine} \& level & \footnotesize Aset audio-visual \\
\midrule
\CPL{03} --- Algoritma dan bahasa & \Utama & \Utama & \Pendukung & \Pendukung \\
\CPL{04} --- Desain, implementasi, evaluasi & \Pendukung & \Pendukung & \Utama & \Utama \\
\bottomrule
\end{tabularx}
\tabsource{Keterangan: \textbf{Utama} = CPMK merupakan perwujudan langsung CPL;
Pendukung = CPMK memperkuat pencapaian CPL.}
\end{table}

\clearpage

\section*{Sub-CPMK}

Setiap CPMK dijabarkan menjadi Sub-CPMK yang lebih spesifik dan terukur.
Setiap Sub-CPMK pada tabel berikut menopang satu bab inti pada buku ini.

\begin{table}[htbp]
\centering
{\setlength{\tabcolsep}{4pt}\footnotesize%
\caption{Sub-CPMK, tingkat taksonomi Bloom, dan CPMK induknya}
\label{tab:subcpmk-buku}
\begin{tabularx}{\textwidth}{@{}C{1.5cm} Y C{1.6cm} C{1.5cm} C{1.7cm}@{}}
\toprule
\headrow \textbf{Kode} & \textbf{Rumusan Sub-CPMK} & \textbf{Bloom} & \textbf{CPMK} & \textbf{CPL} \\
\midrule
1.1 & Menjelaskan struktur \istilah{game loop} (input, \istilah{update},
\istilah{render}) dan peran manajemen waktu dalam permainan. &
C2 & \CPMK{1} & \CPL{03} \\
\addlinespace
1.2 & Merancang dan mengimplementasikan gerak objek berbasis \istilah{delta
time} dan vektor pada Pygame. & C3 & \CPMK{1} & \CPL{03} \\
\addlinespace
1.3 & Menerapkan logika fisika sederhana (gravitasi, lompatan, pantulan)
dalam simulasi 2D. & C3 & \CPMK{1} & \CPL{03} \\
\addlinespace
1.4 & Mendeteksi dan menangani \istilah{collision} menggunakan \code{Rect}
dan/atau \istilah{mask} pada \istilah{sprite}. &
C4 & \CPMK{1} & \CPL{03}, \CPL{04} \\
\addlinespace
2.1 & Merancang perilaku NPC sederhana (misalnya \istilah{chase})
berdasarkan posisi pemain. & C3 & \CPMK{2} & \CPL{03} \\
\addlinespace
2.2 & Mengimplementasikan \istilah{Finite State Machine} (idle, patrol,
chase, attack/flee) untuk NPC. & C3 & \CPMK{2} & \CPL{03} \\
\addlinespace
2.3 & Menerapkan \istilah{pathfinding} berbasis grid dengan algoritma A*
untuk navigasi NPC. & C4 & \CPMK{2} & \CPL{03} \\
\addlinespace
2.4 & Mengevaluasi perilaku NPC hasil kombinasi FSM dan A* terhadap
skenario permainan. & C5 & \CPMK{2} & \CPL{03}, \CPL{04} \\
\addlinespace
3.1 & Menavigasi antarmuka Unity dan menyusun \istilah{project},
\istilah{scene}, serta \istilah{GameObject} dasar. &
C3 & \CPMK{3} & \CPL{04} \\
\addlinespace
3.2 & Membangun level 2D/3D sederhana dengan \istilah{collider},
\istilah{Rigidbody}, dan kontrol pemain. & C4 & \CPMK{3} & \CPL{04} \\
\addlinespace
4.1 & Mengimpor dan mengonfigurasi aset visual (\istilah{sprite}/model,
animasi, UI) ke dalam \istilah{scene}. & C3 & \CPMK{4} & \CPL{04} \\
\addlinespace
4.2 & Mengintegrasikan aset audio dan memicu interaksi audio-visual
berdasarkan peristiwa \istilah{game}. & C4 & \CPMK{4} & \CPL{04} \\
\addlinespace
4.3 & Mengevaluasi prototipe \istilah{game} melalui \istilah{playtesting}
dan memperbaiki integrasi level serta aset. & C5 & \CPMK{4} & \CPL{04} \\
\bottomrule
\end{tabularx}}
\end{table}

\clearpage

\section*{Peta Bab, Materi, dan Minggu Perkuliahan}

Tabel berikut menghubungkan setiap bab pada buku ini dengan materi, Sub-CPMK,
dan minggu pelaksanaan pada RPS. Kolom \istilah{Teknologi} menunjukkan
perkakas utama yang dipakai pada bab tersebut.

\begin{table}[htbp]
\centering
{\setlength{\tabcolsep}{4pt}\footnotesize%
\caption{Peta bab buku ajar terhadap materi, Sub-CPMK, dan minggu RPS}
\label{tab:peta-bab}
\begin{tabularx}{\textwidth}{@{}C{1.2cm} L{4.3cm} C{1.5cm} C{1.8cm} C{1.5cm} Y@{}}
\toprule
\headrow \textbf{Bab} & \textbf{Judul Bab} & \textbf{Materi} & \textbf{Sub-CPMK} & \textbf{Minggu} & \textbf{Teknologi} \\
\midrule
I   & Pendahuluan dan Orientasi Buku & orientasi & --- & 1 & --- \\
II  & Landasan Teori dan Konsep Dasar & fondasi & --- & 1 & --- \\
III & Pengantar Pemrograman Game dan Persiapan Lingkungan & 1 & menuju 1.1 & 1 & Python/Pygame \\
IV  & \istilah{Game Loop}, \istilah{Event Handling}, dan \istilah{Clock} & 2 & 1.1 & 2 & Python/Pygame \\
V   & Gerak Objek, \code{Vector2}, Input, dan \istilah{Delta Time} & 3 & 1.2 & 3 & Python/Pygame \\
VI  & Fisika Sederhana: Gravitasi, Lompatan, dan Pantulan & 4 & 1.3 & 4 & Python/Pygame \\
VII & \istilah{Sprite}, \istilah{Group}, dan \istilah{Collision Detection} & 5 & 1.4 & 5 & Python/Pygame \\
VIII& Pengantar AI NPC dan Perilaku \istilah{Chase} & 6 & 2.1 & 6 & Python/Pygame \\
IX  & \istilah{Finite State Machine} untuk NPC & 7 & 2.2 & 7 & Python/Pygame \\
X   & \istilah{Pathfinding} Grid: dari BFS/Dijkstra ke A* & 8 & 2.3 & 8 & Python/Pygame \\
XI  & Integrasi FSM dan A* (Mini Proyek AI) & 9 & 2.4 & 9 & Python/Pygame \\
---  & \itshape Ujian Tengah Semester & --- & 1.1--2.4 & 10 & --- \\
XII & Pengenalan Unity: \istilah{Project}, \istilah{Scene}, \istilah{GameObject} & 10 & 3.1 & 11 & Unity \\
XIII& Membangun Level, \istilah{Physics}, dan Kontrol Pemain & 11 & 3.2 & 12 & Unity \\
XIV & Impor Aset Visual, Animasi, dan UI Dasar & 12 & 4.1 & 13 & Unity \\
XV  & Integrasi Audio dan \istilah{Trigger} Interaksi & 13 & 4.2 & 14 & Unity \\
XVI & Evaluasi Prototipe, \istilah{Playtesting}, dan Presentasi & 14 & 4.3 & 15 & Unity \\
---  & \itshape Ujian Akhir Semester & --- & 3.1--4.3 & 16 & --- \\
XVII& Evaluasi, Refleksi, dan Integrasi Kompetensi & penutup & 1.1--4.3 & --- & Python \& Unity \\
\bottomrule
\end{tabularx}}
\tabsource{Sumber: disusun dari RPS mata kuliah \NamaMK\ (\KodeMK).}
\end{table}

\begin{catatabox}
Bagian akhir buku memuat empat lampiran pendukung: Lampiran~A rubrik
penilaian lengkap, Lampiran~B contoh tugas dan lembar kerja penilaian,
Lampiran~C glosarium istilah, dan Lampiran~D daftar sumber belajar
tambahan. Peta bobot penilaian terhadap CPMK diuraikan pada Lampiran~A.
\end{catatabox}

\clearpage


\cleardoublepage
\tableofcontents

\cleardoublepage
\listoffigures

\cleardoublepage
\listoftables

% =========================================================
%  MAIN MATTER
%  Bab diaktifkan satu per satu setelah selesai dan diperiksa.
% =========================================================

\mainmatter

\pagestyle{mainmatter}

\subfile{chapters/01/chapter}   % Bab I    Pendahuluan dan Orientasi Buku
\subfile{chapters/02/chapter}   % Bab II   Landasan Teori dan Konsep Dasar
\subfile{chapters/03/chapter}   % Bab III  Materi 1  Pengantar & Setup Python/Pygame
\subfile{chapters/04/chapter}   % Bab IV   Materi 2  Game Loop, Event Handling, Clock
\subfile{chapters/05/chapter}   % Bab V    Materi 3  Gerak Objek, Vector2, Delta Time
\subfile{chapters/06/chapter}   % Bab VI   Materi 4  Fisika Sederhana
\subfile{chapters/07/chapter}   % Bab VII  Materi 5  Sprite, Group, Collision
\subfile{chapters/08/chapter}   % Bab VIII Materi 6  AI NPC dan Perilaku Chase
\subfile{chapters/09/chapter}   % Bab IX   Materi 7  Finite State Machine
\subfile{chapters/10/chapter}   % Bab X    Materi 8  Pathfinding A*
\subfile{chapters/11/chapter}   % Bab XI   Materi 9  Integrasi FSM dan A*
\subfile{chapters/12/chapter}   % Bab XII  Materi 10 Pengenalan Unity
\subfile{chapters/13/chapter}   % Bab XIII Materi 11 Level, Physics, Kontrol Pemain
\subfile{chapters/14/chapter}   % Bab XIV  Materi 12 Aset Visual, Animasi, UI
\subfile{chapters/15/chapter}   % Bab XV   Materi 13 Audio dan Trigger Interaksi
\subfile{chapters/16/chapter}   % Bab XVI  Materi 14 Playtesting dan Presentasi
\subfile{chapters/17/chapter}   % Bab XVII Evaluasi, Refleksi, Integrasi Kompetensi

% =========================================================
%  LAMPIRAN
% =========================================================

\appendix

% appendices/appendix-a-rubrik.tex
\chapter{Rubrik Penilaian Lengkap}
\label{app:rubrik}

Lampiran ini memuat rubrik penilaian dalam bentuk lengkap, yaitu beserta
uraian setiap tingkat capaian. Bab~XVI dan Bab~XVII telah menyajikan bobot dan
kriteria pokoknya; yang ditambahkan di sini adalah perbedaan antara tingkat
yang satu dan tingkat berikutnya, agar penilaian tidak bergantung pada kesan
penilai.

\section{Cara Membaca Rubrik}

Setiap kriteria dinilai pada empat tingkat. Tingkat yang dipakai adalah
tingkat terendah yang \emph{seluruh} syaratnya terpenuhi, bukan tingkat
tertinggi yang sebagian syaratnya terpenuhi. Apabila sebuah pekerjaan memenuhi
tiga dari empat syarat pada tingkat ``Sangat Baik'', ia dinilai ``Baik''.

\begin{table}[htbp]
\centering
\caption{Empat tingkat capaian yang dipakai pada seluruh rubrik}
\label{tab:tingkat-capaian}
{\setlength{\tabcolsep}{4pt}\footnotesize%
\begin{tabularx}{\textwidth}{@{}L{3.4cm} C{2.2cm} Y@{}}
\toprule
\headrow \textbf{Tingkat} & \textbf{Rentang nilai} & \textbf{Ciri umum} \\
\midrule
Sangat Baik & 85--100 & Seluruh syarat terpenuhi, dan pekerjaannya dapat dipakai
sebagai contoh bagi kelompok lain. \\
\addlinespace
Baik & 70--84 & Seluruh syarat pokok terpenuhi, dengan satu kekurangan kecil
yang tidak mengubah hasil. \\
\addlinespace
Cukup & 60--69 & Syarat pokok terpenuhi sebagian; hasilnya berjalan, tetapi
masih menuntut perbaikan sebelum dapat dipakai. \\
\addlinespace
Kurang & $<$ 60 & Syarat pokok tidak terpenuhi; hasilnya tidak berjalan atau
tidak dapat diperiksa. \\
\bottomrule
\end{tabularx}}
\end{table}

\section{Rubrik Mini Proyek AI}

Mini Proyek AI menilai \CPMK{2}, yaitu perilaku NPC. Bobotnya 10 persen dari
nilai akhir, dan dikerjakan berkelompok.

\begin{landscape}
\begin{table}[htbp]
\centering
\caption{Rubrik Mini Proyek AI beserta uraian empat tingkat capaian}
\label{tab:rubrik-mini-lengkap}
{\setlength{\tabcolsep}{4pt}\footnotesize%
\begin{tabularx}{\linewidth}{@{}L{2.8cm} C{1.3cm} Y Y Y Y@{}}
\toprule
\headrow \textbf{Kriteria} & \textbf{Bobot} & \textbf{Sangat Baik} & \textbf{Baik} & \textbf{Cukup} & \textbf{Kurang} \\
\midrule
Implementasi FSM & 25\% &
Keadaan terpisah dan jelas; setiap transisi berjalan menurut jarak dan
\istilah{cooldown}; tidak ada keadaan yang tidak dapat ditinggalkan. &
Transisi berjalan dengan benar; satu keadaan kurang lengkap, misalnya tidak
memiliki jalan kembali. &
Sebagian transisi tidak berjalan; NPC sempat diam tanpa alasan yang jelas. &
Keadaan tidak terpisah; perilaku ditulis sebagai rangkaian \code{if} yang
saling bertumpuk. \\
\addlinespace
Implementasi A* & 25\% &
Jalur benar pada peta berhalangan; memakai antrean prioritas; heuristik
dipilih dengan alasan yang dapat dipertanggungjawabkan. &
Jalur benar dan memakai antrean prioritas; heuristik dipakai tanpa penjelasan. &
Jalur benar pada peta sederhana, tetapi gagal atau lambat pada peta
berhalangan. &
Jalur tidak selalu benar, atau pencarian tidak berakhir pada peta tertentu. \\
\addlinespace
Integrasi FSM dan A* & 20\% &
Pencarian dipicu tepat pada keadaan pengejar; pencarian ulang berkala; kinerja
terjaga dan alasannya dijelaskan. &
Pencarian dipicu pada keadaan yang tepat, tanpa pencarian ulang berkala. &
Pencarian dipicu pada keadaan yang kurang tepat, sehingga jalurnya tertinggal
dari pergerakan pemain. &
Pencarian dijalankan setiap bingkai, sehingga permainan melambat. \\
\addlinespace
Pengujian dan laporan & 20\% &
Sekurang-kurangnya tiga skenario uji; pengamatan objektif dengan angka;
perbaikan ditindaklanjuti dan hasilnya dicatat. &
Tiga skenario uji dan pengamatan objektif; perbaikan belum ditindaklanjuti. &
Skenario uji kurang dari tiga, atau pengamatannya berupa kesan. &
Tidak ada skenario uji maupun laporan. \\
\addlinespace
Presentasi dan kerja sama & 10\% &
Demonstrasi jelas dan tepat waktu; setiap anggota memahami bagiannya dan dapat
menjawab pertanyaan. &
Demonstrasi jelas; satu anggota kurang menguasai bagiannya. &
Demonstrasi bergantung pada satu orang. &
Demonstrasi tidak disiapkan. \\
\bottomrule
\end{tabularx}}
\end{table}
\end{landscape}

\section{Rubrik Proyek Akhir}

Proyek Akhir menilai \CPMK{3} dan \CPMK{4}, dengan bobot 30 persen dari nilai
akhir. Rubrik ini adalah instrumen terbesar pada mata kuliah ini, dan
satu-satunya yang menilai kemampuan menggabungkan bagian-bagian.

\begin{landscape}
\begin{table}[htbp]
\centering
\caption{Rubrik Proyek Akhir beserta uraian empat tingkat capaian}
\label{tab:rubrik-akhir-lengkap}
{\setlength{\tabcolsep}{4pt}\footnotesize%
\begin{tabularx}{\linewidth}{@{}L{2.8cm} C{1.3cm} Y Y Y Y@{}}
\toprule
\headrow \textbf{Kriteria} & \textbf{Bobot} & \textbf{Sangat Baik} & \textbf{Baik} & \textbf{Cukup} & \textbf{Kurang} \\
\midrule
Fungsionalitas mekanik inti dan level & 30\% &
Level selesai dari awal sampai akhir tanpa penghambat; mekanik inti bekerja
sesuai rancangan; kendali pemain responsif; tidak ditemukan
\istilah{soft-lock}. &
Level selesai, dan mekanik intinya bekerja; satu mekanik tambahan belum
berjalan. &
Level dapat dimainkan, tetapi ada bagian yang membingungkan atau satu tempat
yang hanya dapat dicapai dengan susah payah. &
Level tidak dapat diselesaikan, atau mekanik intinya tidak bekerja. \\
\addlinespace
Integrasi aset visual, animasi, dan UI & 25\% &
Aset terpasang rapi dan seragam ukurannya; animasi berganti menurut keadaan;
antarmuka menampilkan angka yang selalu benar. &
Ketiganya terpasang; satu keadaan animasi atau satu angka antarmuka belum
tepat. &
Aset terpasang, tetapi animasi atau antarmuka belum terhubung sepenuhnya pada
keadaan permainan. &
Aset belum terpasang, atau antarmuka menampilkan angka yang tidak sesuai
keadaan. \\
\addlinespace
Integrasi audio dan umpan balik & 15\% &
Musik latar dan efek suara terpisah; setiap efek suara terpicu pada
peristiwanya; tingkat kekuatan beralasan; ada pasangan yang terlihat. &
Efek suara terpicu dengan benar; tingkat kekuatan belum ditata, atau satu
peristiwa belum memperoleh bunyi. &
Sebagian bunyi terpicu, atau musik tertutup oleh efek suara. &
Bunyi belum terpasang, atau terdengar pada peristiwa yang salah. \\
\addlinespace
Kualitas kode dan struktur \istilah{project} & 15\% &
\istilah{Project} tersusun rapi; penamaan konsisten; tidak ada kode mati;
bagian yang berulang disatukan. &
Struktur rapi dan penamaan konsisten; ada sedikit kode mati. &
Struktur dapat dipahami, tetapi penamaannya tidak konsisten atau ada kode yang
tidak terpakai. &
Seluruh kode berada pada satu berkas, atau strukturnya tidak dapat dipahami. \\
\addlinespace
\istilah{Playtest}, laporan, dan presentasi & 15\% &
Lembar \istilah{playtest} terisi oleh dua penguji atau lebih; \istilah{log bug}
lengkap dengan status; laporan memuat keadaan sebelum dan sesudah; demo tepat
waktu. &
Kelengkapan lengkap; satu bagian laporan belum memuat bukti. &
Kelengkapan sebagian; \istilah{log bug} tidak memuat status penanganan. &
Kelengkapan tidak diserahkan, atau demo tidak dilakukan. \\
\bottomrule
\end{tabularx}}
\end{table}
\end{landscape}

\section{Rubrik Portofolio Praktikum}

Tugas praktikum dan portofolio menilai \CPMK{1} sampai \CPMK{4} dengan bobot
25 persen. Penilaiannya dilakukan atas artefak mingguan pada repositori,
bukan atas satu pekerjaan besar.

\begin{table}[htbp]
\centering
\caption{Rubrik portofolio praktikum}
\label{tab:rubrik-portofolio}
{\setlength{\tabcolsep}{4pt}\footnotesize%
\begin{tabularx}{\textwidth}{@{}L{3.4cm} C{1.4cm} Y@{}}
\toprule
\headrow \textbf{Kriteria} & \textbf{Bobot} & \textbf{Deskripsi capaian} \\
\midrule
Kelengkapan artefak & 40\% & Seluruh praktikum terkumpul, masing-masing dengan
kode, keluaran, dan catatan singkat. \\
\addlinespace
Kebenaran hasil & 30\% & Kode berjalan dan menghasilkan keluaran yang sesuai
dengan tujuan praktikum. \\
\addlinespace
Catatan dan penjelasan & 20\% & Setiap artefak disertai penjelasan mengapa
caranya dipilih, bukan hanya apa yang dikerjakan. \\
\addlinespace
Riwayat penyimpanan & 10\% & Pekerjaan disimpan bertahap dengan pesan
penyimpanan yang menerangkan isinya. \\
\bottomrule
\end{tabularx}}
\end{table}

\section{Rubrik Laporan dan Refleksi}

Rubrik ini dipakai bersama Proyek Akhir, dan disajikan lengkap pada
Tabel~\ref{tab:rubrik-laporan-lengkap}.

\begin{table}[htbp]
\centering
\caption{Rubrik laporan pengembangan dan refleksi beserta tingkat capaian}
\label{tab:rubrik-laporan-lengkap}
{\setlength{\tabcolsep}{4pt}\footnotesize%
\begin{tabularx}{\textwidth}{@{}L{3.0cm} C{1.2cm} Y Y@{}}
\toprule
\headrow \textbf{Kriteria} & \textbf{Bobot} & \textbf{Sangat Baik atau Baik} & \textbf{Cukup atau Kurang} \\
\midrule
Kelengkapan dokumen & 20\% &
Memuat ringkasan desain, daftar aset beserta lisensinya, lembar
\istilah{playtest}, dan \istilah{log bug} dengan status penanganan. &
Salah satu kelengkapan tidak diserahkan, atau daftar aset tidak memuat
lisensinya. \\
\addlinespace
Laporan perbaikan & 30\% &
Setiap perbaikan disertai keadaan sebelum, keadaan sesudah, dan buktinya. &
Perbaikan diceritakan tanpa bukti, atau hanya sebagian perbaikan yang
dilaporkan. \\
\addlinespace
Kesesuaian GDD dan \istilah{build} & 20\% &
Perbedaan dicantumkan dengan jujur menurut tiga keadaan pada
Tabel~\ref{tab:gdd-build}, beserta alasannya. &
Perbedaan tidak dicantumkan, atau janji dihapus dari GDD agar tampak sesuai. \\
\addlinespace
Refleksi teknologi & 30\% &
Setiap pernyataan memuat keadaan konkret, tindakan, dan hasil; ketiga tinjauan
dijawab terpisah. &
Refleksi berupa kesan umum, atau hanya sebagian tinjauan yang dijawab. \\
\bottomrule
\end{tabularx}}
\end{table}

\section{Ketentuan Kelulusan}

Kelulusan mata kuliah ini ditentukan oleh tiga ketentuan yang berlaku
sekaligus: nilai akhir sekurang-kurangnya 60, kehadiran sekurang-kurangnya
75 persen, dan penyerahan seluruh kelengkapan Proyek Akhir. Bobot setiap
komponen, matriks sumbangannya kepada CPMK, dan konversi nilainya disajikan
pada Bagian~\ref{sec:bobot-nilai}.
          % Lampiran A
% appendices/appendix-b-tugas-portofolio.tex
\chapter{Contoh Tugas, Portofolio, dan Lembar Kerja}
\label{app:lembar-kerja}

Lampiran ini memuat bentuk-bentuk yang dipakai sepanjang perkuliahan dalam
keadaan siap dicetak. Berkas aslinya tersedia pada repositori mata kuliah
dalam bentuk yang dapat disunting; yang disajikan di sini adalah bentuk dan
isinya, agar Anda mengetahui apa yang diharapkan sebelum mengerjakannya.

\section{Format Laporan Praktikum}

Setiap praktikum pada Bab~III sampai Bab~XV menghasilkan satu berkas laporan.
Laporan itu tidak perlu panjang; satu halaman sudah cukup, asalkan memuat
keempat bagian pada Tabel~\ref{tab:format-laporan}.

\begin{table}[htbp]
\centering
\caption{Format laporan praktikum satu halaman}
\label{tab:format-laporan}
{\setlength{\tabcolsep}{4pt}\footnotesize%
\begin{tabularx}{\textwidth}{@{}L{3.0cm} Y@{}}
\toprule
\headrow \textbf{Bagian} & \textbf{Isi, dan mengapa diperlukan} \\
\midrule
Tujuan & Satu atau dua kalimat yang menyebut Sub-CPMK dan hasil yang hendak
dicapai. Bagian ini memudahkan Anda memeriksa sendiri apakah praktikumnya
berhasil. \\
\addlinespace
Yang dikerjakan & Uraian singkat langkah yang diambil, termasuk perubahan yang
tidak ada pada petunjuk. Bagian ini yang membedakan laporan hasil kerja dari
salinan petunjuk. \\
\addlinespace
Keluaran & Tangkapan layar atau cuplikan keluaran program, beserta angkanya
apabila ada. Keluaran yang disertai angka dapat diperiksa ulang. \\
\addlinespace
Catatan & Masalah yang ditemui, dugaan penyebabnya, dan cara mengatasinya.
Bagian ini yang paling berguna bagi Anda pada praktikum berikutnya. \\
\bottomrule
\end{tabularx}}
\end{table}

\section{Format Portofolio}

Portofolio adalah kumpulan seluruh artefak praktikum pada satu repositori.
Susunan yang dianjurkan diperlihatkan pada Tabel~\ref{tab:format-portofolio}.
Susunan ini meniru susunan sumber buku ini, sehingga mahasiswa yang mengikuti
seluruh praktikum akan memperoleh kebiasaan yang sama.

\begin{table}[htbp]
\centering
\caption{Susunan direktori portofolio}
\label{tab:format-portofolio}
{\setlength{\tabcolsep}{4pt}\footnotesize%
\begin{tabularx}{\textwidth}{@{}L{4.4cm} Y@{}}
\toprule
\headrow \textbf{Direktori} & \textbf{Isi} \\
\midrule
\code{code/python/} & Seluruh berkas Python dari Bab~III sampai Bab~XI, dengan
nama yang menyebut babnya. \\
\addlinespace
\code{code/unity/} & Seluruh berkas C\# dari Bab~XII sampai Bab~XVI. \\
\addlinespace
\code{laporan/} & Laporan praktikum satu halaman untuk setiap pertemuan. \\
\addlinespace
\code{aset/} & Aset yang dipakai, beserta berkas \code{SUMBER.md} yang memuat
sumber dan lisensi setiap aset. \\
\addlinespace
\code{proyek/} & Proyek akhir: \istilah{build}, kode sumber, laporan
pengembangan, lembar \istilah{playtest}, dan \istilah{log bug}. \\
\bottomrule
\end{tabularx}}
\end{table}

Aset berukuran besar sebaiknya tidak disimpan langsung di dalam repositori.
Simpanlah pada tempat penyimpanan berkas, lalu catat alamatnya pada berkas
\code{SUMBER.md}.

\section{Lembar \istilah{Playtest} Siap Cetak}

Lembar pada Tabel~\ref{tab:lembar-playtest-cetak} dipakai pada sesi
\istilah{playtest} antar kelompok. Isilah satu lembar untuk setiap penguji.

\begin{table}[htbp]
\centering
\caption{Lembar \istilah{playtest} siap cetak}
\label{tab:lembar-playtest-cetak}
{\setlength{\tabcolsep}{4pt}\footnotesize%
\begin{tabularx}{\textwidth}{@{}L{4.0cm} Y@{}}
\toprule
\headrow \textbf{Butir} & \textbf{Isian} \\
\midrule
Nama penguji dan kelompok & \\[1.3em]
\midrule
Pengalaman bermain sejenis & \\[1.3em]
\midrule
Nama \istilah{build} dan tanggalnya & \\[1.3em]
\midrule
Waktu mulai dan waktu selesai & \\[1.3em]
\midrule
Tujuan sesi (bagian yang diragukan) & \\[1.3em]
\midrule
Titik berhenti lebih dari 10 detik & \\[2.0em]
\midrule
Keadaan yang tidak seharusnya terjadi & \\[2.0em]
\midrule
Pendapat penguji, dicatat apa adanya & \\[2.0em]
\midrule
Catatan pengamat sendiri & \\[2.0em]
\bottomrule
\end{tabularx}}
\end{table}

\section{Log \istilah{Bug} Siap Cetak}

Log pada Tabel~\ref{tab:log-bug-cetak} memuat kolom yang sama dengan
Tabel~\ref{tab:log-bug} pada Bab~XVI, dalam bentuk yang dapat diisi langsung.

\begin{table}[htbp]
\centering
\caption{Log \istilah{bug} siap cetak}
\label{tab:log-bug-cetak}
{\setlength{\tabcolsep}{4pt}\footnotesize%
\begin{tabularx}{\textwidth}{@{}C{0.7cm} C{1.6cm} Y C{1.3cm} C{1.4cm}@{}}
\toprule
\headrow \textbf{No} & \textbf{Jenis} & \textbf{Urutan yang menyebabkannya, dan hasil yang salah} & \textbf{Berat} & \textbf{Status} \\
\midrule
1 & & & & \\[2.2em]
\midrule
2 & & & & \\[2.2em]
\midrule
3 & & & & \\[2.2em]
\midrule
4 & & & & \\[2.2em]
\midrule
5 & & & & \\[2.2em]
\bottomrule
\end{tabularx}}
\end{table}

Jenis temuan yang dipakai adalah mekanik, integrasi, tampilan, dan
\istilah{soft-lock}. Beratnya adalah berat, sedang, atau ringan, menurut
Tabel~\ref{tab:berat-temuan}. Statusnya adalah belum, sedang dikerjakan, atau
selesai beserta tanggalnya.

\section{Daftar Periksa Penyerahan Proyek Akhir}

Serahkan seluruh butir pada Tabel~\ref{tab:penyerahan} dalam satu
\istilah{build} dan satu repositori. Berkas yang hilang pada hari penyerahan
tidak dapat disusulkan tanpa persetujuan dosen pengampu.

\begin{table}[htbp]
\centering
\caption{Daftar periksa penyerahan Proyek Akhir}
\label{tab:penyerahan}
{\setlength{\tabcolsep}{4pt}\footnotesize%
\begin{tabularx}{\textwidth}{@{}C{0.8cm} L{6.2cm} Y@{}}
\toprule
\headrow \textbf{No} & \textbf{Kelengkapan} & \textbf{Keterangan} \\
\midrule
1 & \istilah{Build} yang dapat dijalankan & Diuji pada komputer lain, bukan
pada komputer pembuatnya. \\
\addlinespace
2 & Kode sumber \istilah{project} & Seluruh berkas yang diperlukan untuk
membangun ulang. \\
\addlinespace
3 & Ringkasan desain (GDD) & Beserta catatan perubahan yang disengaja. \\
\addlinespace
4 & Laporan pengembangan & Memuat keadaan sebelum dan sesudah setiap
perbaikan. \\
\addlinespace
5 & Lembar \istilah{playtest} terisi & Sekurang-kurangnya dua penguji dari
kelompok lain. \\
\addlinespace
6 & \istilah{Log bug} & Beserta status penanganan setiap temuan. \\
\addlinespace
7 & Daftar aset dan lisensinya & Memuat sumber setiap aset yang bukan karya
sendiri. \\
\addlinespace
8 & Lembar refleksi akhir semester & Menurut Bagian~\ref{sec:refleksi-akhir}. \\
\bottomrule
\end{tabularx}}
\end{table}

\section{Contoh Tugas yang Dapat Diberikan Dosen}

Tabel~\ref{tab:contoh-tugas} memuat contoh tugas tambahan yang dapat dipakai
apabila suatu CPMK belum mencapai ambang batas. Setiap tugas menunjuk pada
bagian buku yang menjadi sumbernya, sehingga mahasiswa dapat mengerjakannya
tanpa bahan tambahan.

\begin{table}[htbp]
\centering
\caption{Contoh tugas tambahan dan bagian buku yang menjadi sumbernya}
\label{tab:contoh-tugas}
{\setlength{\tabcolsep}{4pt}\footnotesize%
\begin{tabularx}{\textwidth}{@{}L{2.0cm} L{4.6cm} Y@{}}
\toprule
\headrow \textbf{CPMK} & \textbf{Bentuk tugas} & \textbf{Sumber} \\
\midrule
\CPMK{1} & Mengukur laju bingkai pada tiga nilai pembatas, lalu menjelaskan
hubungannya dengan beban kerja & Bab~IV, dan \code{ch04_ukur_laju.py} \\
\addlinespace
\CPMK{1} & Membandingkan gerak berbasis jumlah bingkai dengan gerak berbasis
\istilah{delta time} pada dua ukuran jendela & Bab~IV dan Bab~V \\
\addlinespace
\CPMK{2} & Menelusuri urutan keadaan NPC pada satu skenario, dan menunjukkan
transisi yang tidak sah & Bab~IX, dan \code{ch09_fsm.py} \\
\addlinespace
\CPMK{2} & Membandingkan banyak petak yang diperiksa oleh dua heuristik pada
peta yang sama & Bab~X, dan \code{ch10_astar.py} \\
\addlinespace
\CPMK{3} & Menjelaskan gejala fisika pada sebuah \istilah{scene}, dan
menentukan komponen penyebabnya & Bab~XIII \\
\addlinespace
\CPMK{4} & Menyusun daftar peristiwa permainan beserta bunyi yang seharusnya
memicunya & Bab~XV \\
\addlinespace
\CPMK{4} & Menjalankan \code{ch16_uji_mandiri.py} atas level sendiri, lalu
memperbaiki tempat yang dilaporkan & Bab~XVI \\
\bottomrule
\end{tabularx}}
\end{table}
% Lampiran B
% appendices/appendix-c-glosarium.tex
\chapter{Glosarium}
\label{app:glosarium}

Glosarium ini memuat istilah yang dipakai pada seluruh buku, disusun menurut
abjad istilah Indonesianya. Kolom kedua memuat padanan Inggrisnya, sebab
sebagian besar dokumentasi dan pesan kesalahan yang akan Anda temui memakai
padanan itu.

\begin{longtable}{@{}L{3.9cm} L{3.6cm} L{7.2cm}@{}}
\caption{Glosarium istilah Indonesia--Inggris}
\label{tab:glosarium} \\
\toprule
\headrow \textbf{Istilah} & \textbf{Padanan Inggris} & \textbf{Keterangan singkat} \\
\midrule
\endfirsthead
\toprule
\headrow \textbf{Istilah} & \textbf{Padanan Inggris} & \textbf{Keterangan singkat} \\
\midrule
\endhead
\bottomrule
\endlastfoot

Algoritma pencarian jalur & \istilah{Pathfinding algorithm} &
Algoritma yang mencari rangkaian petak dari satu titik ke titik lain, misalnya
BFS, Dijkstra, dan A*. \\ \addlinespace

Animasi & \istilah{Animation} &
Rangkaian gambar yang diputar berurutan sehingga tampak bergerak. \\
\addlinespace

Antarmuka pengguna & \istilah{User interface} (UI) &
Bagian layar yang bukan bagian dari dunia permainan, misalnya skor, nyawa, dan
tombol. \\ \addlinespace

Antrean prioritas & \istilah{Priority queue} &
Struktur data yang mengeluarkan isinya menurut prioritas, bukan menurut urutan
masuk. A* memakainya untuk memilih petak berikutnya. \\ \addlinespace

Aset & \istilah{Asset} &
Berkas yang dipakai permainan: gambar, bunyi, huruf, dan animasi. \\
\addlinespace

Batang (badan) & \istilah{Body} &
Benda yang diperlakukan sebagai satu kesatuan gerak, dengan posisi, kecepatan,
dan bentuk. \\ \addlinespace

Batas benda & \istilah{Collider} &
Bentuk yang dipakai mesin untuk memeriksa tabrakan, tidak selalu sama dengan
bentuk yang terlihat pemain. \\ \addlinespace

Bingkai & \istilah{Frame} &
Satu kali gambar lengkap ditampilkan. Pada 60 bingkai per detik, satu bingkai
berlangsung sekitar 16,7 milidetik. \\ \addlinespace

Bunyi latar & \istilah{Background music} (BGM) &
Bunyi panjang yang berjalan terus selama permainan, dan tidak dipanggil
menurut kejadian. \\ \addlinespace

Bug & \istilah{Bug} &
Keadaan yang tidak seharusnya terjadi, dan dapat diulang dengan urutan
tindakan tertentu. \\ \addlinespace

Build & \istilah{Build} &
Berkas permainan yang sudah dibungkus sehingga dapat dijalankan pada komputer
lain. \\ \addlinespace

Celah & \istilah{Gap} &
Jarak kosong yang harus dilompati pemain di antara dua pijakan. \\
\addlinespace

Delta time & \istilah{Delta time} &
Selisih waktu antara dua bingkai, dipakai agar gerak tidak bergantung pada
kecepatan komputer. \\ \addlinespace

Desibel & \istilah{Decibel} (dB) &
Satuan kekuatan bunyi yang naik secara berlipat. Turun 6 dB berarti kekuatannya
menjadi kira-kira separuh. \\ \addlinespace

Deteksi tabrakan & \istilah{Collision detection} &
Pemeriksaan apakah dua benda saling bersentuhan pada satu bingkai. \\
\addlinespace

Efek suara & \istilah{Sound effect} (SFX) &
Bunyi pendek yang dipanggil menurut peristiwa, dan dapat berbunyi bersamaan
dengan efek suara lain. \\ \addlinespace

Gaya berat & \istilah{Gravity} &
Percepatan tetap yang menarik benda ke bawah. Pada Bab~VI dipakai 1800 piksel
per detik kuadrat. \\ \addlinespace

Graf & \istilah{Graph} &
Kumpulan simpul dan hubungan di antaranya. Peta yang diubah menjadi pijakan
dan lompatan adalah sebuah graf. \\ \addlinespace

Heuristik & \istilah{Heuristic} &
Perkiraan jarak ke tujuan yang dipakai A* untuk mengarahkan pencarian. \\
\addlinespace

Hierarki & \istilah{Hierarchy} &
Susunan induk dan anak pada \istilah{scene}, misalnya \istilah{Canvas} yang
membawahi teks dan ikon. \\ \addlinespace

Impor & \istilah{Import} &
Memasukkan berkas aset ke dalam \istilah{project}, sehingga berkas itu diolah
dan disimpan dalam bentuk yang siap dipakai. \\ \addlinespace

Indikator & \istilah{Indicator} &
Perilaku yang dapat diamati, yang dipakai untuk menyatakan bahwa sebuah
Sub-CPMK telah dikuasai. \\ \addlinespace

Jangkauan & \istilah{Reach} &
Jarak terjauh yang dapat ditempuh pemain dalam satu lompatan, dihitung dari
kecepatan lompatan dan gravitasi. \\ \addlinespace

Kecerdasan buatan & \istilah{Artificial intelligence} (AI) &
Pada buku ini, aturan yang menentukan perilaku NPC, bukan pembelajaran mesin. \\
\addlinespace

Kecepatan bingkai & \istilah{Frame rate} &
Banyaknya bingkai yang ditampilkan dalam satu detik, dinyatakan dalam
\istilah{fps}. \\ \addlinespace

Keadaan & \istilah{State} &
Keadaan yang ditempati sebuah benda atau sistem pada satu saat, misalnya diam,
berlari, atau melompat. \\ \addlinespace

Kendali pemain & \istilah{Player controller} &
Skrip yang membaca masukan pemain dan mengubahnya menjadi gerak benda. \\
\addlinespace

Kode sumber & \istilah{Source code} &
Berkas teks berisi program sebelum dijalankan atau dibangun. \\
\addlinespace

Komponen & \istilah{Component} &
Bagian yang ditempelkan pada \istilah{GameObject} untuk memberi kemampuan,
misalnya \istilah{Rigidbody} atau \istilah{Animator}. \\ \addlinespace

Lembar \istilah{playtest} & \istilah{Playtest sheet} &
Lembar yang dipakai penguji untuk mencatat waktu, tempat, dan keadaan yang
tidak seharusnya terjadi. \\ \addlinespace

Level & \istilah{Level} &
Satu ruang permainan beserta seluruh isinya, dari titik awal sampai tujuan. \\
\addlinespace

Log \istilah{bug} & \istilah{Bug log} &
Daftar temuan beserta urutan yang menyebabkannya, beratnya, dan status
penanganannya. \\ \addlinespace

Masukan & \istilah{Input} &
Apa pun yang dikirim pemain ke program: tombol, tetikus, atau sentuhan. \\
\addlinespace

Mekanik & \istilah{Mechanic} &
Aturan yang menentukan cara pemain memengaruhi permainan, misalnya melompat
dan mengumpulkan benda. \\ \addlinespace

Mesin keadaan & \istilah{State machine} (FSM) &
Susunan keadaan beserta aturan perpindahan di antaranya. \\
\addlinespace

Mesin permainan & \istilah{Game engine} &
Perangkat lunak yang mengurus tampilan, fisika, bunyi, dan aset, misalnya
Unity. \\ \addlinespace

Nada & \istilah{Pitch} &
Tinggi rendahnya bunyi. Menggeser nada sedikit membuat bunyi yang berulang
tidak terdengar seperti mesin. \\ \addlinespace

NPC & \istilah{Non-player character} &
Benda yang bertindak menurut aturan program, bukan menurut masukan pemain. \\
\addlinespace

Objek permainan & \istilah{GameObject} &
Satuan dasar pada Unity: benda yang memiliki letak, dan memperoleh kemampuan
melalui komponen. \\ \addlinespace

Partikel & \istilah{Particle} &
Gambar kecil yang muncul, bergerak, lalu menghilang, dipakai sebagai pasangan
yang terlihat bagi bunyi. \\ \addlinespace

Peristiwa & \istilah{Event} &
Kejadian yang dibaca satu kali, bukan keadaan yang berlangsung, misalnya
tombol ditekan atau benda disentuh. \\ \addlinespace

Petak & \istilah{Tile} &
Satuan kotak pada peta dua dimensi. Jangkauan lompatan pada buku ini dinyatakan
dalam petak. \\ \addlinespace

Pijakan & \istilah{Platform} &
Petak kosong yang di bawahnya berupa dinding, sehingga pemain dapat berdiri di
sana. \\ \addlinespace

\istilah{Pixels Per Unit} & \istilah{Pixels per unit} &
Banyaknya piksel gambar yang mewakili satu satuan pada \istilah{scene}.
Menentukan ukuran gambar di dalam level. \\ \addlinespace

\istilah{Playtest} & \istilah{Playtest} &
Sesi ketika orang lain memainkan prototipe Anda, sementara Anda mengamati dan
mencatat. \\ \addlinespace

Portofolio & \istilah{Portfolio} &
Kumpulan artefak pekerjaan beserta keterangannya, disusun sebagai bukti
kemampuan. \\ \addlinespace

Prototipe & \istilah{Prototype} &
Versi awal yang sudah dapat dimainkan, dipakai untuk menguji rancangan sebelum
dikerjakan sepenuhnya. \\ \addlinespace

\istilah{Rigidbody} & \istilah{Rigid body} &
Komponen Unity yang menyerahkan gerak benda kepada mesin fisika, sehingga
gravitasi dan tabrakan dihitung oleh mesin. \\ \addlinespace

Simulasi & \istilah{Simulation} &
Menjalankan permainan tanpa pemain, berulang kali, untuk memperoleh gambaran
tentang sebagian keadaannya. \\ \addlinespace

Skrip & \istilah{Script} &
Berkas berisi kode yang ditempelkan pada \istilah{GameObject} sebagai komponen
pada Unity. \\ \addlinespace

\istilah{Soft-lock} & \istilah{Soft-lock} &
Keadaan ketika pemain tidak dapat melanjutkan dan tidak dapat kembali,
walaupun programnya masih berjalan dan tidak mencetak kesalahan. \\
\addlinespace

\istilah{Spatial Blend} & \istilah{Spatial blend} &
Pengaturan pada \code{AudioSource} yang menentukan seberapa besar letak bunyi
berpengaruh: nol untuk musik latar, satu untuk bunyi di dalam level. \\
\addlinespace

\istilah{Sprite} & \istilah{Sprite} &
Gambar dua dimensi yang ditempatkan di dalam \istilah{scene}. \\
\addlinespace

\istilah{Sprite sheet} & \istilah{Sprite sheet} &
Satu berkas gambar yang memuat beberapa bingkai, dipotong menjadi bingkai
terpisah. \\ \addlinespace

Sub-CPMK & \istilah{Sub-learning outcome} &
Turunan CPMK pada satu pokok bahasan, beserta indikator ketercapaiannya. \\
\addlinespace

Tabrakan & \istilah{Collision} &
Persentuhan dua benda yang mengubah geraknya, berbeda dari
\istilah{trigger} yang hanya mendeteksi sentuhan. \\ \addlinespace

Tembus & \istilah{Tunneling} &
Keadaan ketika benda berpindah melewati benda lain dalam satu langkah, karena
langkahnya terlalu besar. \\ \addlinespace

Tingkat (audio) & \istilah{Mixer group} &
Kelompok pada \istilah{AudioMixer} yang mengatur kekuatan sekumpulan bunyi
sekaligus. \\ \addlinespace

Transform & \istilah{Transform} &
Komponen yang menyimpan letak, rotasi, dan ukuran sebuah
\istilah{GameObject}. \\ \addlinespace

\istilah{Trigger} & \istilah{Trigger} &
Bentuk batas benda yang mendeteksi sentuhan tanpa menghentikan gerak, dipakai
untuk pintu, benda yang dapat diambil, dan daerah tujuan. \\
\addlinespace

Umpan balik & \istilah{Feedback} &
Tanggapan yang diterima pemain atas tindakannya, melalui tampilan, bunyi, atau
gerak. \\ \addlinespace

Vektor & \istilah{Vector} &
Besaran yang memiliki besar dan arah, dipakai untuk menyatakan posisi dan
kecepatan. \\ \addlinespace

Waktu bingkai & \istilah{Frame time} &
Lama satu bingkai berlangsung. Kebalikan dari kecepatan bingkai. \\
\addlinespace

\end{longtable}
       % Lampiran C
% appendices/appendix-d-sumber-belajar.tex
\chapter{Sumber Belajar}
\label{app:sumber}

Lampiran ini menjelaskan kegunaan setiap rujukan pada Daftar Pustaka, agar
Anda dapat memilih bacaan menurut keperluan. Uraian lengkap setiap rujukan
tidak diulang di sini; yang ditambahkan hanyalah keterangan tentang bagian
mana yang berguna dan untuk bab mana.

\section{Buku Teks dan Dokumentasi Resmi}

Rujukan pada Tabel~\ref{tab:sumber-utama} adalah bacaan pokok. Ketiganya
tersedia bebas, dan ketiganya dipakai berulang kali sepanjang buku ini.

\begin{table}[htbp]
\centering
\caption{Buku teks dan dokumentasi resmi beserta kegunaannya}
\label{tab:sumber-utama}
{\setlength{\tabcolsep}{4pt}\footnotesize%
\begin{tabularx}{\textwidth}{@{}C{1.1cm} L{4.6cm} Y@{}}
\toprule
\headrow \textbf{Kode} & \textbf{Bacaan} & \textbf{Kegunaan pada buku ini} \\
\midrule
\cite{sweigart-mg} & Sweigart, \istilah{Making Games with Python \& Pygame} &
Pembahasan permainan lengkap dengan Pygame, dari \istilah{game loop} sampai
bunyi. Berguna sebagai contoh utuh setelah Bab~XI. \\
\addlinespace
\cite{sweigart-invent} & Sweigart, \istilah{Invent Your Own Computer Games with
Python} & Dasar Python dan permainan berbasis teks. Berguna sebelum Bab~III
apabila dasar Python Anda belum kuat. \\
\addlinespace
\cite{pygame-docs} & Dokumentasi resmi Pygame &
Rujukan perintah yang dipakai pada Bab~III--XI. Periksa di sini ketika nama
atau argumen sebuah fungsi tidak lagi sesuai dengan contoh pada buku. \\
\addlinespace
\cite{pygame-sprite-mask} & Dokumentasi \code{pygame.mask} &
Pembahasan tabrakan berdasarkan bentuk gambar, yang tidak dibahas pada
Bab~VII tetapi berguna untuk permainan dengan bentuk tidak beraturan. \\
\addlinespace
\cite{unity-leveldesign} & Unity Technologies, \istilah{Introduction to Level
Design in Game Development and in Unity} &
Rancangan level, penyampaian keterangan kepada pemain, dan pengujian. Menjadi
sandaran Bab~XII--XVI. \\
\addlinespace
\cite{unity-learn} & Unity Learn &
Kumpulan tutorial resmi Unity. Berguna untuk melatih ulang langkah pada
Bab~XII--XV sebelum mengerjakan proyek. \\
\addlinespace
\cite{unity-manual-import} & Unity Manual: \istilah{Importing Assets} dan
\istilah{Audio} &
Rujukan pengaturan impor dan bunyi pada Bab~XIV--XV. \\
\bottomrule
\end{tabularx}}
\end{table}

\section{Tulisan Teknis Mendalam}

Rujukan pada Tabel~\ref{tab:sumber-teknis} membahas satu topik secara
mendalam. Sebagian besar tidak perlu dibaca seluruhnya; bacalah bagian yang
bersangkutan ketika Anda ingin memahami alasan di balik sebuah pilihan.

\begin{table}[htbp]
\centering
\caption{Tulisan teknis mendalam beserta bab yang berkaitan}
\label{tab:sumber-teknis}
{\setlength{\tabcolsep}{4pt}\footnotesize%
\begin{tabularx}{\textwidth}{@{}C{1.1cm} L{4.6cm} Y@{}}
\toprule
\headrow \textbf{Kode} & \textbf{Bacaan} & \textbf{Kegunaan pada buku ini} \\
\midrule
\cite{nystrom-gameloop} & Nystrom, ``Fix Your Timestep'' dan pembahasan
\istilah{game loop} & Alasan di balik \istilah{game loop} pada Bab~IV, termasuk
pemisahan antara pembaruan tetap dan pembaruan tampilan. \\
\addlinespace
\cite{fiedler-timestep} & Fiedler, \istilah{Timestep} &
Perbandingan beberapa cara menjalankan waktu pada permainan. Berguna setelah
Bab~IV apabila gerak permainan Anda tidak stabil. \\
\addlinespace
\cite{catto-integration} & Catto, \istilah{Integration Basics} &
Penjelasan Euler dan Verlet untuk fisika sederhana. Sandaran Bab~VI, termasuk
alasannya mengenai galat yang mengecil ketika langkah waktunya diperkecil. \\
\addlinespace
\cite{santos-pathfinding} & Santos, \istilah{Pathfinding} &
Penjelasan BFS, Dijkstra, dan A* beserta ilustrasinya. Sandaran Bab~X. \\
\addlinespace
\cite{santos-fsm} & Santos, \istilah{Finite State Machines} &
Penjelasan mesin keadaan beserta contohnya. Sandaran Bab~IX. \\
\addlinespace
\cite{bendersky-pygame-astar} & Bendersky, A* pada Pygame &
Contoh penerapan A* pada permainan dua dimensi. Berguna sebagai pembanding
\code{ch10_astar.py}. \\
\addlinespace
\cite{munir-astar} & Munir, \istilah{Pencarian Jalur dengan A*} &
Pembahasan sifat heuristik dan pengaruhnya pada hasil pencarian. Berguna untuk
menjawab latihan pada Bab~X. \\
\addlinespace
\cite{itb-astar-navigasi} & ITB, ``Implementasi Algoritma A* dalam Navigasi
Karakter'' & Contoh penerapan A* pada karakter permainan. \\
\addlinespace
\cite{stmikdharmapala-astar} & STMIK Dharmapala, penerapan A* pada permainan &
Contoh penerapan pada permainan buatan Indonesia. \\
\addlinespace
\cite{itb-collision} & ITB, ``Collision Detection in Precision Parry-Based
Combat'' & Pembahasan deteksi tabrakan yang lebih rinci daripada Bab~VII. \\
\addlinespace
\cite{pygame-chimp} & Pygame Community, \istilah{Chimp Line-by-Line} &
Contoh permainan Pygame yang dibahas baris demi baris. Berguna setelah
Bab~VII untuk melihat susunan program secara utuh. \\
\addlinespace
\cite{reintech-game-ai} & Reintech, ``How to Create a Video Game AI with
Python'' & Contoh perilaku NPC dalam Python. Berguna sebagai pembanding
Bab~VIII. \\
\bottomrule
\end{tabularx}}
\end{table}

\section{Rencana Pembelajaran dan Kurikulum Pembanding}

Rujukan pada Tabel~\ref{tab:sumber-rps} dipakai untuk menyusun peta capaian
pada Bab~XVII, dan berguna bagi Anda yang ingin mengetahui bagaimana mata
kuliah serupa disusun di tempat lain.

\begin{table}[htbp]
\centering
\caption{Rencana pembelajaran dan kurikulum pembanding}
\label{tab:sumber-rps}
{\setlength{\tabcolsep}{4pt}\footnotesize%
\begin{tabularx}{\textwidth}{@{}C{1.1cm} L{4.6cm} Y@{}}
\toprule
\headrow \textbf{Kode} & \textbf{Bacaan} & \textbf{Kegunaan pada buku ini} \\
\midrule
\cite{its-rps-game} & ITS, \istilah{RPS Desain dan Pemrograman Game} &
Pembanding susunan capaian dan pokok bahasan. \\
\addlinespace
\cite{unima-rps-game} & UNIMA, \istilah{RPS Pengembangan Game} &
Pembanding bentuk penilaian dan presentasi proyek. \\
\addlinespace
\cite{unesa-rps-game} & UNESA, \istilah{RPS Pemrograman Game} &
Pembanding urutan materi antara blok algoritma dan blok mesin permainan. \\
\addlinespace
\cite{polibatam-kurikulum} & Politeknik Negeri Batam, \istilah{Kurikulum
Program Studi Teknologi Permainan} & Pembanding daftar periksa capaian dan
penyusunan daftar periksa pada Bab~XVI. \\
\addlinespace
\cite{univrab-pemrograman-game} & Universitas Abdurrab, \istilah{Deskripsi Mata
Kuliah Pemrograman Game} & Pembanding cakupan materi dan prasyarat. \\
\bottomrule
\end{tabularx}}
\end{table}

\section{Karya Akademik Indonesia}

Rujukan pada Tabel~\ref{tab:sumber-karya} adalah skripsi, jurnal, dan laporan
pengabdian dari perguruan tinggi Indonesia. Kegunaannya berbeda dari buku
teks: karya seperti ini memperlihatkan bagaimana sebuah permainan sungguh
dibangun dan dilaporkan, termasuk kekurangannya.

\begin{table}[htbp]
\centering
\caption{Karya akademik Indonesia beserta kegunaannya}
\label{tab:sumber-karya}
{\setlength{\tabcolsep}{4pt}\footnotesize%
\begin{tabularx}{\textwidth}{@{}C{1.1cm} L{4.6cm} Y@{}}
\toprule
\headrow \textbf{Kode} & \textbf{Bacaan} & \textbf{Kegunaan pada buku ini} \\
\midrule
\cite{uma-pygame} & Universitas Medan Area, \istilah{Panduan Membuat Game 2D
Sederhana dengan Pygame} & Contoh pengantar yang ringkas untuk Bab~III--IV. \\
\addlinespace
\cite{ub-pygame} & Universitas Brawijaya, ``Game Development with Python'' &
Contoh penyusunan permainan sederhana dan penjelasannya. \\
\addlinespace
\cite{dinus-pygame} & Universitas Dian Nuswantoro, ``Pemrograman dan Game
dengan Python/Pygame'' & Contoh pelatihan Pygame bagi pemula. \\
\addlinespace
\cite{stmikjayakarta-pygame} & STMIK Jayakarta, ``Penerapan Library Pygame
dalam Game RPG The Adventure'' & Contoh permainan RPG dengan Pygame, termasuk
pengurusan bunyi. \\
\addlinespace
\cite{itn-astar-fsm} & ITN Malang, ``Penerapan A* dan FSM pada Game Lost
Civilization'' & Contoh penggabungan A* dan FSM, sebagaimana Bab~XI. \\
\addlinespace
\cite{umri-astar-behaviour} & UMRI, ``Penerapan Algoritma A* dan Behaviour
Trees'' & Pembanding pendekatan lain untuk perilaku NPC. \\
\addlinespace
\cite{uinmalang-fsm-fuzzy} & UIN Malang, \istilah{Implementasi FSM dan Logika
Fuzzy} & Pembanding pendekatan yang lebih halus daripada FSM pada Bab~IX. \\
\addlinespace
\cite{uny-katelu} & UNY, \istilah{Pengembangan Permainan Edukasi KATELU dengan
Unity 3D} & Contoh proyek Unity lengkap beserta laporannya. \\
\addlinespace
\cite{undip-media} & UNDIP, ``Media Pembelajaran Budaya Indonesia Menggunakan
Unity'' & Contoh pencantuman aset dan sumber pada karya Unity. \\
\addlinespace
\cite{untan-edukasi} & UNTAN, ``Game Edukasi Budaya Indonesia'' &
Pembanding rancangan permainan edukasi. \\
\addlinespace
\cite{rokania-unity} & STKIP Rokania, \istilah{Pelatihan Dasar Pemrograman Game
dengan Unity} & Contoh pelatihan Unity bagi pemula, dan pembanding refleksi
pemilihan teknologi pada Bab~XVI. \\
\bottomrule
\end{tabularx}}
\end{table}

\section{Saran Urutan Membaca}

Buku ini disusun agar dapat dibaca berurutan. Apabila Anda memerlukan bacaan
tambahan, urutan berikut biasanya paling berguna.

\begin{enumerate}
  \item \textbf{Sebelum Bab~III}: \cite{sweigart-invent} untuk dasar Python.
  \item \textbf{Setelah Bab~IV}: \cite{nystrom-gameloop} dan
        \cite{fiedler-timestep} untuk memahami mengapa waktunya dijalankan
        dengan cara itu.
  \item \textbf{Setelah Bab~VI}: \cite{catto-integration} untuk galat
        integrasi dan cara memperkecilnya.
  \item \textbf{Setelah Bab~IX dan Bab~X}: \cite{santos-fsm} dan
        \cite{santos-pathfinding} untuk pembahasan yang lebih rinci.
  \item \textbf{Setelah Bab~XI}: \cite{sweigart-mg} untuk melihat satu
        permainan Pygame yang lengkap.
  \item \textbf{Setelah Bab~XVI}: \cite{unity-leveldesign} bagian pengujian,
        lalu salah satu karya akademik pada Tabel~\ref{tab:sumber-karya}
        sebagai contoh proyek yang dilaporkan.
\end{enumerate}
  % Lampiran D

% =========================================================
%  BACK MATTER
% =========================================================

\backmatter

% =========================================================
%  references.tex
%  Daftar pustaka buku ajar.
%  Nomor rujukan [R1]--[R35] sengaja dibuat sama dengan
%  rps/sections/10-referensi.tex agar tabel pemetaan minggu
%  dan rujukan pada RPS tetap berlaku pada buku ajar.
%  \cite{kunci} menghasilkan [Rn] sesuai nomor tersebut.
% =========================================================

\begin{sloppypar}
% Lingkungan thebibliography pada kelas book mencetak judul tanpa
% menambahkannya ke daftar isi, sehingga baris berikut diperlukan.
\addcontentsline{toc}{chapter}{\bibname}
\begin{thebibliography}{R35}
\setlength{\itemsep}{5pt}

% ---------- Referensi utama (dipakai mahasiswa) ----------
\bibitem[R1]{sweigart-mg}
  A. Sweigart, \istilah{Making Games with Python \& Pygame}, 2012.
  \url{https://inventwithpython.com/makinggames.pdf}

\bibitem[R2]{sweigart-invent}
  A. Sweigart, \istilah{Invent Your Own Computer Games with Python},
  edisi ke-4, 2018.
  \url{https://inventwithpython.com/inventwithpython.pdf}

\bibitem[R3]{pygame-docs}
  Pygame Community, \istilah{Pygame Documentation --- Quick Start dan
  Referensi API}. \url{https://www.pygame.org/docs/}

\bibitem[R4]{pygame-sprite-mask}
  Pygame Community, \istilah{Pygame Documentation --- Sprite} dan
  \istilah{Mask}. \url{https://www.pygame.org/docs/ref/sprite.html} dan
  \url{https://www.pygame.org/docs/ref/mask.html}

\bibitem[R5]{nystrom-gameloop}
  R. Nystrom, \istilah{Game Programming Patterns} --- bab
  \istilah{Game Loop}.
  \url{https://gameprogrammingpatterns.com/game-loop.html}

\bibitem[R6]{fiedler-timestep}
  G. Fiedler, ``Fix Your Timestep!'', \istilah{Gaffer On Games}, 2004.
  \url{https://gafferongames.com/post/fix_your_timestep/}

\bibitem[R7]{catto-integration}
  E. Catto, \istilah{Numerical Integration}, GDC 2009.
  \url{https://box2d.org/files/ErinCatto_NumericalIntegration_GDC2009.pdf}

\bibitem[R8]{santos-pathfinding}
  E. dos Santos, \istilah{Game AI: Pathfinding}, bahan kuliah, 2020.
  \url{https://edirlei.com/aulas/game-ai-2020/GAME_AI_Lecture_02_Pathfinding_2020.pdf}

\bibitem[R9]{santos-fsm}
  E. dos Santos, \istilah{Game AI: Finite State Machines}, bahan kuliah,
  2020.
  \url{https://edirlei.com/aulas/game-ai-2020/GAME_AI_Lecture_03_Finite_State_Machines_2020.pdf}

\bibitem[R10]{bendersky-pygame-astar}
  E. Bendersky, ``Writing a Game in Python with Pygame, Part III:
  Pathfinding'', 2009.
  \url{https://eli.thegreenplace.net/2009/01/09/writing-a-game-in-python-with-pygame-part-iii}

\bibitem[R11]{munir-astar}
  R. Munir, \istilah{Penentuan Rute, Bagian 2: Algoritma A*}, bahan
  kuliah Institut Teknologi Bandung.
  \url{https://informatika.stei.itb.ac.id/~rinaldi.munir/Stmik/2020-2021/Route-Planning-Bagian2-2021.pdf}

\bibitem[R12]{unity-leveldesign}
  Unity Technologies, \istilah{Introduction to Level Design in Game
  Development and in Unity}, e-book resmi.
  \url{https://assets2.brandfolder.io/bf-boulder-prod/khmctn58vjqkr42wmm3mqvq9/v/1131033155/original/UNITY_Ebook_Introduction%20to%20level%20design%20in%20game%20development%20and%20in%20Unity_V2.pdf}

\bibitem[R13]{unity-learn}
  Unity Technologies, \istilah{Unity Learn --- Unity Essentials} dan
  \istilah{Junior Programmer: Create with Code}.
  \url{https://learn.unity.com/}

\bibitem[R14]{unity-manual-import}
  Unity Technologies, \istilah{Unity Manual --- Importing Assets} dan
  \istilah{Audio}.
  \url{https://docs.unity3d.com/Manual/ImportingAssets.html}

% ---------- Referensi pendukung: RPS dan materi perguruan tinggi ----------
\bibitem[R15]{its-rps-game}
  Institut Teknologi Sepuluh Nopember, \istilah{RPS Desain dan
  Pemrograman Game}, Program Studi Teknik Komputer.
  \url{https://www.its.ac.id/komputer/wp-content/uploads/sites/28/2018/03/Desain-dan-pemograman-Game.pdf}

\bibitem[R16]{unima-rps-game}
  Universitas Negeri Manado, \istilah{RPS Pengembangan Game}, Program
  Studi Teknik Informatika, 2025/2026.
  \url{https://ti.unima.ac.id/rps/view/11}

\bibitem[R17]{unesa-rps-game}
  Universitas Negeri Surabaya, \istilah{RPS Pemrograman Game}, S1
  Pendidikan Teknologi Informasi.
  \url{https://sidia.unesa.ac.id/rps-pdf/s1-pendidikan-teknologi-informasi/pemrograman-game.pdf}

\bibitem[R18]{polibatam-kurikulum}
  Politeknik Negeri Batam, \istilah{Kurikulum Program Studi Teknologi
  Permainan}, 2024.
  \url{https://www.polibatam.ac.id/wp-content/uploads/2024/06/Kurikulum-Prodi-TP.pdf}

\bibitem[R19]{univrab-pemrograman-game}
  Universitas Abdurrab, \istilah{Deskripsi Mata Kuliah Pemrograman Game},
  SPADA. \url{http://spada.univrab.ac.id/course/info.php?id=1303&lang=en}

\bibitem[R20]{uma-pygame}
  Universitas Medan Area, \istilah{Panduan Membuat Game 2D Sederhana
  dengan Pygame}, 2025.
  \url{https://p2dpm.uma.ac.id/2025/10/15/panduan-membuat-game-2d-sederhana-dengan-pygame/}

\bibitem[R21]{ub-pygame}
  Universitas Brawijaya, ``Game Development with Python: Let's Build a
  Simple Game!'', \istilah{Gramaswara}.
  \url{https://gramaswara.ub.ac.id/index.php/gramaswara/article/download/185/128}

\bibitem[R22]{dinus-pygame}
  Universitas Dian Nuswantoro, ``Pemrograman dan Game dengan
  Python/Pygame'', \istilah{Abdimasku}.
  \url{https://abdimasku.lppm.dinus.ac.id/index.php/jurnalabdimasku/article/download/2965/735}

\bibitem[R23]{itb-collision}
  Institut Teknologi Bandung, ``Collision Detection in Precision
  Parry-Based Combat'', makalah IF2123, 2024.
  \url{https://informatika.stei.itb.ac.id/~rinaldi.munir/AljabarGeometri/2024-2025/Makalah/Makalah-IF2123-Algeo-2024%20(31).pdf}

\bibitem[R24]{stmikjayakarta-pygame}
  STMIK Jayakarta, ``Penerapan Library Pygame dalam Game RPG The
  Adventure'', \istilah{Jurnal JMI Jayakarta}.
  \url{https://journal.stmikjayakarta.ac.id/index.php/JMIJayakarta/article/download/1013/677/}

\bibitem[R25]{itn-astar-fsm}
  Institut Teknologi Nasional Malang, ``Penerapan A* Pathfinding dan
  Finite State Machine pada Game Lost Civilization'', \istilah{JATI}.
  \url{https://ejournal.itn.ac.id/index.php/jati/article/download/5402/3593}

\bibitem[R26]{umri-astar-behaviour}
  Universitas Muhammadiyah Riau, ``Penerapan Algoritma A* dan Behaviour
  Trees untuk NPC'', \istilah{COSCITECH}.
  \url{https://ejurnal.umri.ac.id/index.php/coscitech/article/download/5418/2453/17612}

\bibitem[R27]{uinmalang-fsm-fuzzy}
  UIN Maulana Malik Ibrahim Malang, \istilah{Implementasi FSM dan Logika
  Fuzzy untuk Perilaku NPC}, skripsi.
  \url{http://etheses.uin-malang.ac.id/8204/1/11650053.pdf}

\bibitem[R28]{uny-katelu}
  Universitas Negeri Yogyakarta, \istilah{Pengembangan Permainan Edukasi
  KATELU dengan Unity 3D}, skripsi.
  \url{https://eprints.uny.ac.id/52540/1/TAS_AFIF%20FATCHUR%20RAHMAN_12520241062.pdf}

\bibitem[R29]{undip-media}
  Universitas Diponegoro, ``Media Pembelajaran Budaya Indonesia
  Menggunakan Unity'', \istilah{Jurnal Teknologi dan Sistem Komputer}.
  \url{https://jtsiskom.undip.ac.id/article/download/12667/12286}

\bibitem[R30]{untan-edukasi}
  Universitas Tanjungpura, ``Game Edukasi Budaya Indonesia'',
  \istilah{JCSKom}.
  \url{https://jurnal.untan.ac.id/index.php/jcskommipa/article/download/49393/75676592053}

\bibitem[R31]{rokania-unity}
  STKIP Rokania, \istilah{Pelatihan Dasar Pemrograman Game dengan Unity},
  \istilah{Jurnal MNR}.
  \url{https://e-jurnal.rokania.ac.id/index.php/jmnr/article/download/39/28/172}

\bibitem[R32]{pygame-chimp}
  Pygame Community, \istilah{Pygame Tutorial: Line-by-Line Chimp}.
  \url{https://www.pygame.org/docs/tut/ChimpLineByLine.html}

\bibitem[R33]{reintech-game-ai}
  Reintech, ``How to Create a Video Game AI with Python''.
  \url{https://reintech.io/blog/how-to-create-a-video-game-ai-with-python}

\bibitem[R34]{itb-astar-navigasi}
  Institut Teknologi Bandung, ``Implementasi Algoritma A* dalam Navigasi
  Karakter'', makalah IF2120, 2023.
  \url{https://informatika.stei.itb.ac.id/~rinaldi.munir/Matdis/2022-2023/Makalah2022/Makalah-Matdis-2022%20(86).pdf}

\bibitem[R35]{stmikdharmapala-astar}
  STMIK Dharmapala Riau, ``Rancang Bangun Pathfinding A* pada NPC
  Menggunakan Unity Engine'', \istilah{Jurnal JIKB}.
  \url{https://ojs.stmikdharmapalariau.ac.id/index.php/jikb/article/download/831/574}

\end{thebibliography}
\end{sloppypar}


\end{document}
